\documentclass[11pt, a4paper]{article}

\usepackage{amsmath, amsthm, amssymb}
\usepackage{mathtools}
\usepackage[margin=2.5cm]{geometry}
\usepackage{tcolorbox}
\usepackage{hyperref}
\usepackage{microtype}
\numberwithin{equation}{section}

\tcbuselibrary{theorems, skins, breakable}

\newtcbtheorem[number within=section]{theorem}{Theorem}{
  enhanced, breakable,
  colback=blue!4!white, colframe=blue!50!black,
  fonttitle=\bfseries, separator sign={.}
}{thm}

\newtcbtheorem[number within=section]{lemma}{Lemma}{
  enhanced, breakable,
  colback=green!4!white, colframe=green!40!black,
  fonttitle=\bfseries, separator sign={.}
}{lem}

\newtcbtheorem[number within=section]{corollary}{Corollary}{
  enhanced, breakable,
  colback=orange!5!white, colframe=orange!55!black,
  fonttitle=\bfseries, separator sign={.}
}{cor}

\newtcbtheorem[number within=section]{remark}{Remark}{
  enhanced, breakable,
  colback=gray!6!white, colframe=gray!50!black,
  fonttitle=\bfseries, separator sign={.}
}{rem}

\newcommand{\norm}[1]{\left\|#1\right\|}
\newcommand{\abs}[1]{\left|#1\right|}
\newcommand{\tX}{\widetilde{X}}
\newcommand{\tD}{\widetilde{D}}
\newcommand{\tT}{\widetilde{T}}
\newcommand{\tU}{\widetilde{U}}
\newcommand{\tV}{\widetilde{V}}
\newcommand{\tSig}{\widetilde{\Sigma}}
\newcommand{\tW}{\widetilde{W}}
\newcommand{\R}{\mathbb{R}}
\newcommand{\C}{\mathbb{C}}

\title{\textbf{Enclosure of Eigenvalues via Approximate Diagonalization}\\[0.4em]
\large A note in the spirit of Miyajima and Bauer--Fike}
\author{}
\date{}

\begin{document}
\maketitle

\section{The bridge: Miyajima certification as a pseudospectral engine}

\subsection{What this note is about}

The $\varepsilon$-pseudospectrum of $A \in \C^{n\times n}$,
\[
  \sigma_\varepsilon(A)
    = \bigl\{ z \in \C : \norm{(A - zI)^{-1}} \ge \varepsilon^{-1} \bigr\}
    = \bigl\{ z \in \C : \sigma_{\min}(A - zI) \le \varepsilon \bigr\},
\]
is a region, not a finite set, and \emph{rigorously enclosing} it means
certifying, over a continuum of test points $z$, both an \emph{inclusion}
(where the resolvent is provably large) and an \emph{exclusion} (where it is
provably small).  The computational bottleneck is the inner loop: a rigorous
bound on $\norm{(A - zI)^{-1}}$ --- equivalently a certified lower bound on
$\sigma_{\min}(A - zI)$ --- evaluated at every $z$ of a sweep.  A region of even
modest resolution requires this certificate at thousands of points, so its
per-point cost, not its existence, decides feasibility.

This note develops an inner-loop certificate built on \emph{Miyajima's
enclosure machinery}.  Miyajima's 2010 theorem was designed to enclose
eigenvalues of a generalized eigenproblem by measuring two residuals $R_1, R_2$
of a computed eigendecomposition, never forming an explicit inverse and folding
all conditioning into a single residual norm $\norm{R_2}$.  We repurpose that
mechanism: the same residuals yield, at each $z$, a rigorous resolvent-norm
bound, and hence a pseudospectral certificate.  The point is efficiency --- the
residual structure makes the per-$z$ cost $O(n^2)$ after a single $O(n^3)$
setup, with all $z$-independent data computed once.

\subsection{The bridge theorem}

The following statement is the conceptual anchor of the note; every later
section is a sharpening, a specialization, or a cost analysis of it.

\begin{theorem}{Miyajima certification yields a pseudospectral certificate}{bridge}
  Let $A \in \C^{n\times n}$, and let $\tX$ (approximate eigenvectors or Schur
  basis), $M$ (diagonal $\tD$ or triangular $\tT$), and $Y \approx \tX^{-1}$ be
  computed once.  Form the $z$-independent products $B = Y\tX$ (so $R_2 = B -
  I$) and $C = YA\tX$, and at each test point the affine residual $R_1(z) = Y(A
  - zI)\tX = C - zB$.  Suppose $\norm{R_2} < 1$.  Then for every $z$:
  \begin{enumerate}
    \item \textup{(Exclusion / resolvent bound.)} If the certified gap
      $g(z) > 0$ produced by the constructions below is positive, then $z
      \notin \sigma(A)$ and
      \[
        \norm{(A - zI)^{-1}} \;\le\; \frac{\kappa(\tX)}{g(z)},
      \]
      so $z \notin \sigma_\varepsilon(A)$ whenever $g(z) > \varepsilon\,
      \kappa(\tX)$.  All ingredients of $g(z)$ and $\kappa(\tX)$ are computed
      from $\tX, Y, R_1(z), R_2$ without forming any inverse.
    \item \textup{(Inclusion.)} Conversely, a certified \emph{lower} bound on
      $\norm{(A - zI)^{-1}}$ (equivalently a certified upper bound on
      $\sigma_{\min}(A - zI)$, supplied by the cluster machinery of Section~3)
      places $z$ inside $\sigma_\varepsilon(A)$.
  \end{enumerate}
  Consequently the certified set $\{ z : g(z) > \varepsilon\,\kappa(\tX) \}$ is a
  rigorous \emph{exclusion region} for $\sigma_\varepsilon(A)$, and its
  complement, refined by the inclusion certificate, encloses the pseudospectrum.
\end{theorem}

\noindent
The theorem is deliberately schematic: the symbol $g(z)$ stands for whichever
certified gap one computes --- the uniform gap of Theorem~\ref{thm:diag}, the
recentred per-row gap of Theorem~\ref{thm:resolvent-row}, the interpolated
$\ell^2$ gap, or the block-assembled gap of Theorem~\ref{thm:assembly} --- and
$\kappa(\tX)$ for the matching conditioning factor.  The rest of the note makes
$g(z)$ progressively sharper and cheaper, and Section~\ref{sec:diag} begins with
the simplest instance.

\subsection{Why Miyajima, and what is new}

Three features of the Miyajima mechanism make it the right engine, and their
combination for pseudospectral enclosure is the contribution of this note.
\begin{itemize}
  \item \emph{No inverse, conditioning hidden in one residual.}  The bound is
    driven by $\norm{R_2} = \norm{Y\tX - I}$; a good $Y \approx \tX^{-1}$ makes
    it small, and no separate $\norm{A}$ or condition-number factor appears in
    the \emph{enclosure}.  Conditioning re-enters only in the resolvent
    \emph{magnitude}, through $\kappa(\tX)$, where it is genuine.
  \item \emph{$z$-independent setup.}  The $z$-independent data are formed once:
    $\tX, M, Y$, the products $B \coloneqq Y\tX$ (hence $R_2 = B - I$) and $C
    \coloneqq YA\tX$, and the reduced solve $\tilde S$.  The per-$z$ residual is
    then \emph{affine} in $z$,
    \[
      R_1(z) = Y(A - zI)\tX = YA\tX - z\,Y\tX = C - z B,
    \]
    an $O(n^2)$ update (a scalar multiply and a subtraction), never a matrix
    product.  This affine structure is what makes sweeping a two-dimensional
    region affordable; the triangular layer below is cheaper still, needing
    nothing per $z$ but the centres.
  \item \emph{Triangular structure keeps $\ell^1, \ell^\infty$ trivial; $\ell^2$
    by interpolation.}  The hard object --- a certified $\sigma_{\min}$ of a
    triangular matrix in $\ell^2$ --- is never computed directly except on small
    clusters; elsewhere the $\ell^1$/$\ell^\infty$ diagonal-dominance gaps and
    Riesz--Thorin suffice.
\end{itemize}
Rigorous pseudospectrum bounding is an active area (Trefethen--Embree for the
classical theory; Colbrook--Hansen and the SCI program, and ResDMD, for the
rigorous and data-driven turns), but those approaches are largely $\ell^2$ /
Hilbert-space and either self-adjoint-flavoured or data-driven.  The present
route is operator-norm based, exploits triangular/Schur structure, and uses
Miyajima certification as its inner loop --- a combination that, to our
knowledge, has not been used for pseudospectral enclosure.  Eigenvalue-cluster
enclosure (Rump 2022, Rump--Lange 2023) and verified matrix decompositions
(Rump--Ogita 2024) supply related certified building blocks but target
points, not the resolvent-norm \emph{surface} that a pseudospectrum requires.
Later sections quantify the efficiency claim by comparing candidate inner-loop
certificates on shared examples.

\section{Approximate diagonalization}
\label{sec:diag}

\subsection{Setup}

Let $A \in \C^{n\times n}$.  Suppose a numerical eigensolver has returned an
approximate eigenvector matrix $\tX \in \C^{n\times n}$ and a diagonal matrix
$\tD = \operatorname{diag}(\tilde d_1,\ldots,\tilde d_n)$ of approximate
eigenvalues, so that $A\tX \approx \tX\tD$.  We do not assume $\tX$ is exactly
invertible; nonsingularity will instead be \emph{verified} as a by-product of
the enclosure.

The construction follows Miyajima's enclosure theorem, a verified form of
Bauer--Fike.  Its efficiency comes from introducing one auxiliary matrix and
measuring two residuals, rather than forming inverses or transpose products.
Let $Y \in \C^{n\times n}$ be arbitrary; in practice $Y$ is computed as an
approximate inverse $Y \approx \tX^{-1}$.  Define the two residuals
\[
  R_1 \coloneqq Y(A\tX - \tX\tD), \qquad
  R_2 \coloneqq Y\tX - I.
\]
Both are single matrix products of computed data and are enclosed rigorously in
interval or ball arithmetic.  We use the $\ell^\infty$ operator norm
$\norm{M}_\infty = \max_i \sum_j \abs{M_{ij}}$ throughout.

\subsection{The enclosure theorem}

\begin{theorem}{Miyajima enclosure (standard eigenproblem)}{diag}
  With $R_1, R_2$ as above, suppose $\norm{R_2}_\infty < 1$.  Then $\tX$ and $Y$
  are nonsingular, and every eigenvalue $\lambda$ of $A$ satisfies
  \[
    \min_{1\le i\le n}\abs{\lambda - \tilde d_i} \;\le\; \varepsilon,
    \qquad
    \varepsilon \coloneqq \frac{\norm{R_1}_\infty}{1 - \norm{R_2}_\infty}.
  \]
  Equivalently, $\sigma(A) \subseteq \bigcup_{i=1}^n \{ z \in \C :
  \abs{z - \tilde d_i} \le \varepsilon \}$.
\end{theorem}

\begin{proof}
Since $\norm{R_2}_\infty < 1$, the Neumann series shows $I + R_2 = Y\tX$ is
nonsingular, hence so are $Y$ and $\tX$; moreover $\norm{(I+R_2)^{-1}}_\infty
\le (1 - \norm{R_2}_\infty)^{-1}$.

Let $\lambda$ be an eigenvalue of $A$.  If $\lambda = \tilde d_i$ for some $i$
the bound is trivial, so assume $\tD - \lambda I$ is nonsingular.  Because $A -
\lambda I$ is singular while $\tD - \lambda I$ and $\tX$ are not, consider the
identity
\[
  (\tD - \lambda I)^{-1}\tX^{-1}(A - \lambda I)\tX
    = I - (\tD - \lambda I)^{-1}\tX^{-1}(\tX\tD - A\tX),
\]
obtained by writing $\tX^{-1}(A-\lambda I)\tX = (\tD - \lambda I) - \tX^{-1}(\tX\tD
- A\tX)$ and multiplying by $(\tD - \lambda I)^{-1}$.  The left-hand side is
singular (as $A - \lambda I$ is), so the right-hand side $I - M$ is singular,
where $M \coloneqq (\tD - \lambda I)^{-1}\tX^{-1}(A\tX - \tX\tD)$.  By the
contrapositive of the Neumann bound, singularity of $I - M$ forces
$\norm{M}_\infty \ge 1$.  Therefore
\[
  1 \le \norm{M}_\infty
    \le \norm{(\tD - \lambda I)^{-1}}_\infty\,\norm{\tX^{-1}(A\tX - \tX\tD)}_\infty
    = \frac{\norm{\tX^{-1}(A\tX - \tX\tD)}_\infty}{\min_i \abs{\lambda - \tilde d_i}},
\]
using that $\tD - \lambda I$ is diagonal, so its inverse has $\ell^\infty$ norm
$1/\min_i\abs{\lambda - \tilde d_i}$.  Hence
\[
  \min_i\abs{\lambda - \tilde d_i}
    \le \norm{\tX^{-1}(A\tX - \tX\tD)}_\infty.
\]
Finally eliminate the unknown inverse $\tX^{-1}$ in favour of the computed $Y$:
since $\tX^{-1} = (Y\tX)^{-1}Y = (I + R_2)^{-1}Y$,
\[
  \tX^{-1}(A\tX - \tX\tD) = (I+R_2)^{-1}Y(A\tX - \tX\tD) = (I + R_2)^{-1}R_1,
\]
so that
\[
  \min_i\abs{\lambda - \tilde d_i}
    \le \norm{(I+R_2)^{-1}}_\infty\,\norm{R_1}_\infty
    \le \frac{\norm{R_1}_\infty}{1 - \norm{R_2}_\infty} = \varepsilon. \qedhere
\]
\end{proof}

\subsection{Individual radii}

Theorem~\ref{thm:diag} gives a single uniform radius $\varepsilon$.  A sharper,
Gerschgorin-style statement assigns each centre $\tilde d_i$ its own radius by
keeping the row structure of $(I+R_2)^{-1}R_1$ instead of collapsing to a norm.

\begin{corollary}{Individual (Gerschgorin-type) radii}{individual}
  Let $S \coloneqq (I + R_2)^{-1} R_1$, and let $\rho_i \coloneqq \sum_{j}
  \abs{S_{ij}}$ be the $i$-th absolute row sum of $S$.  Under the hypothesis
  $\norm{R_2}_\infty < 1$ of Theorem~\textup{\ref{thm:diag}}, there is for each
  eigenvalue $\lambda$ of $A$ an index $i^*$ with
  \[
    \abs{\lambda - \tilde d_{i^*}} \le \rho_{i^*}.
  \]
  Consequently $\sigma(A) \subseteq \bigcup_{i=1}^n \{ z : \abs{z - \tilde d_i}
  \le \rho_i \}$.
\end{corollary}

\begin{proof}
Let $v$ be an eigenvector, $Av = \lambda v$, and set $u \coloneqq \tX^{-1}v \neq
0$ (valid since $\tX$ is nonsingular).  From $A\tX = \tX\tD + (A\tX - \tX\tD)$
and $Av = \lambda v$,
\[
  \lambda v = A\tX u = \tX\tD u + (A\tX - \tX\tD)u
  \;\Longrightarrow\;
  \tX(\lambda I - \tD)u = (A\tX - \tX\tD)u,
\]
and applying $\tX^{-1}$, then writing $\tX^{-1}(A\tX-\tX\tD) = (I+R_2)^{-1}R_1 =
S$,
\[
  (\lambda I - \tD)u = S u,
  \qquad\text{i.e.}\qquad
  (\lambda - \tilde d_i)\,u_i = (Su)_i \ \text{ for each } i.
\]
Choose $i^*$ with $\abs{u_{i^*}} = \norm{u}_\infty > 0$ and divide:
\[
  \abs{\lambda - \tilde d_{i^*}}
    = \frac{\abs{(Su)_{i^*}}}{\norm{u}_\infty}
    \le \frac{\sum_j \abs{S_{i^*j}}\abs{u_j}}{\norm{u}_\infty}
    \le \sum_j \abs{S_{i^*j}} = \rho_{i^*}.
\]
Since $\rho_{i^*} \le \max_i \rho_i = \norm{S}_\infty \le \varepsilon$, the
corollary refines Theorem~\ref{thm:diag}.
\end{proof}

In computed terms, using \eqref{eq:rho-tilde}, the discs may be taken as
$B(\tilde d_i,\ \tilde\rho_i + \beta_\infty)$ with $\tilde\rho_i = \sum_j
\abs{\tilde S_{ij}}$ and $\beta_\infty = \norm{(I+R_2)\tilde S - R_1}_\infty /
(1 - \norm{R_2}_\infty)$ --- all quantities computed, with no inverse formed.

\subsection{The computed reduced residual and its slack}

The matrix $S \coloneqq (I + R_2)^{-1} R_1$ governs every bound below, but it
contains an inverse and so is not a bare product like $R_1, R_2$.  Following the
same residual philosophy that controls the first inverse through $R_2$, we never
form $(I+R_2)^{-1}$.  Instead we take \emph{any} computed approximation $\tilde
S \approx S$ --- e.g.\ a floating-point solve of $(I + R_2)\tilde S = R_1$ ---
and pay only for its residual.  Define
\[
  T \coloneqq (I + R_2)\tilde S - R_1, \qquad
  \tau_p \coloneqq \norm{T}_p, \qquad
  \beta_p \coloneqq \frac{\tau_p}{1 - \norm{R_2}_p} \quad (p \in \{1,\infty\}).
\]

\begin{lemma}{Residual control of $S$}{Sresidual}
  If $\norm{R_2}_p < 1$, then $\norm{S - \tilde S}_p \le \beta_p$.
\end{lemma}

\begin{proof}
$S - \tilde S = (I+R_2)^{-1}R_1 - \tilde S = -(I+R_2)^{-1}\bigl((I+R_2)\tilde S
- R_1\bigr) = -(I+R_2)^{-1}T$, so $\norm{S - \tilde S}_p \le
\norm{(I+R_2)^{-1}}_p\,\norm{T}_p \le \tau_p/(1 - \norm{R_2}_p) = \beta_p$.
\end{proof}

Thus every appearance of $S$ may be replaced by the computed $\tilde S$ at the
cost of a single uniform slack $\beta_p$.  Two consequences are used repeatedly.
First, for the row/column sums: writing $\tilde\rho_i \coloneqq \sum_j
\abs{\tilde S_{ij}}$,
\begin{equation}\label{eq:rho-tilde}
  \rho_i = \sum_j \abs{S_{ij}} \le \tilde\rho_i + \beta_\infty,
\end{equation}
since the $i$-th row of $S - \tilde S$ has $\ell^1$ norm at most $\beta_\infty$.
Second, the recentred gap of Theorem~\ref{thm:resolvent-row} degrades by at most
$\beta_\infty$, because the centre shift $\abs{S_{ii} - \tilde S_{ii}}$ and the
radius perturbation $\sum_{j\ne i}\abs{S_{ij} - \tilde S_{ij}}$ together sum to
the full row deviation $\sum_j \abs{S_{ij} - \tilde S_{ij}} \le \beta_\infty$.

\begin{remark}{Crude vs.\ genuine $\tilde S$}{crude}
  Taking $\tilde S = 0$ gives $T = -R_1$ and $\beta_\infty = \norm{R_1}_\infty /
  (1 - \norm{R_2}_\infty)$: the bound reverts to charging the entire $R_1$.  A
  genuine floating-point solve makes $\tau_p$ the (tiny) solve residual, so
  $\beta_p$ is negligible and the computed $\tilde S$ carries the real
  information --- including its diagonal $\tilde S_{ii}$, which supplies the free
  recentring $\tilde d_i \to \tilde d_i - \tilde S_{ii}$.  Crucially $\tilde S$
  and $T$ are independent of $\lambda$, so this one-time cost amortises over an
  entire sweep in $\lambda$.
\end{remark}



The same identity that yields the enclosure also bounds the resolvent
$(A - \lambda I)^{-1}$ for $\lambda$ outside the discs.  Unlike the enclosure,
the resolvent bound \emph{does} carry the conditioning of $\tX$ --- this is
unavoidable, as the resolvent genuinely measures non-normality --- but every
factor remains rigorously computable from $\tX, Y, R_1, R_2$.

For $p \in \{1, \infty\}$ write $\norm{\cdot}_p$ for the induced operator norm
and $S \coloneqq (I + R_2)^{-1} R_1$.  Recall from the proof of
Theorem~\ref{thm:diag} that
\[
  \tX^{-1}(A - \lambda I)\tX = (\tD - \lambda I)(I - M),
  \qquad M = (\tD - \lambda I)^{-1} S,
\]
so that, when $I - M$ is invertible,
\begin{equation}\label{eq:resolvent-factored}
  (A - \lambda I)^{-1} = \tX\,(I - M)^{-1}(\tD - \lambda I)^{-1}\tX^{-1}.
\end{equation}

\begin{theorem}{Resolvent bound, $\ell^p$ for $p \in \{1,\infty\}$}{resolvent}
  Suppose $\norm{R_2}_p < 1$ and let $g_p(\lambda) \coloneqq \min_i \abs{\lambda
  - \tilde d_i} - \norm{S}_p$.  If $g_p(\lambda) > 0$ then $\lambda \notin
  \sigma(A)$ and
  \[
    \norm{(A - \lambda I)^{-1}}_p
      \;\le\; \frac{\norm{\tX}_p\,\norm{\tX^{-1}}_p}{g_p(\lambda)}
      \;\le\; \frac{\norm{\tX}_p\,\norm{Y}_p}
                   {\bigl(1 - \norm{R_2}_p\bigr)\,g_p(\lambda)}.
  \]
\end{theorem}

\begin{proof}
Since $\tD - \lambda I$ is diagonal, $\norm{(\tD - \lambda I)^{-1}}_p =
1/\min_i\abs{\lambda - \tilde d_i}$ for both $p = 1$ and $p = \infty$ (the
maximum absolute row sum and column sum of a diagonal matrix coincide).  Hence
$\norm{M}_p \le \norm{S}_p / \min_i\abs{\lambda - \tilde d_i} < 1$ exactly when
$g_p(\lambda) > 0$, in which case $I - M$ is nonsingular (so $A - \lambda I$ is,
by \eqref{eq:resolvent-factored}) with $\norm{(I-M)^{-1}}_p \le (1 -
\norm{M}_p)^{-1}$.  Taking $\norm{\cdot}_p$ in \eqref{eq:resolvent-factored},
\[
  \norm{(A - \lambda I)^{-1}}_p
    \le \norm{\tX}_p\,\norm{\tX^{-1}}_p \cdot
        \frac{1}{1 - \norm{M}_p}\cdot\frac{1}{\min_i\abs{\lambda - \tilde d_i}}
    \le \frac{\norm{\tX}_p\,\norm{\tX^{-1}}_p}
             {\min_i\abs{\lambda - \tilde d_i} - \norm{S}_p},
\]
where the last step uses $\frac{1}{1-\norm{M}_p}\cdot\frac{1}{m} \le
\frac{1}{m - \norm{S}_p}$ with $m = \min_i\abs{\lambda - \tilde d_i}$ and
$\norm{M}_p \le \norm{S}_p/m$.  Finally $\norm{\tX^{-1}}_p =
\norm{(I+R_2)^{-1}Y}_p \le \norm{Y}_p/(1 - \norm{R_2}_p)$.
\end{proof}

The two bounds interpolate.  Writing $\kappa_p(\tX) = \norm{\tX}_p
\norm{\tX^{-1}}_p$ for the $\ell^p$ condition number, Riesz--Thorin gives the
$\ell^2$ resolvent bound at no extra cost.

\begin{theorem}{Interpolated resolvent bound, $\ell^2$}{resolvent2}
  Suppose $\norm{R_2}_1 < 1$ and $\norm{R_2}_\infty < 1$, and that both
  $g_1(\lambda) > 0$ and $g_\infty(\lambda) > 0$.  Then $\lambda \notin
  \sigma(A)$ and
  \[
    \norm{(A - \lambda I)^{-1}}_2
      \;\le\; \sqrt{\,\norm{(A - \lambda I)^{-1}}_1\;
                     \norm{(A - \lambda I)^{-1}}_\infty\,}
      \;\le\; \sqrt{\frac{\kappa_1(\tX)\,\kappa_\infty(\tX)}
                         {g_1(\lambda)\,g_\infty(\lambda)}}.
  \]
  More generally, for any $\theta \in [0,1]$ with $\tfrac1{p} = 1-\theta$,
  $\tfrac1{q} = \theta$ \textup{(}so $\tfrac12 = \tfrac{1-\theta}{1} +
  \tfrac{\theta}{\infty}$ at $\theta = \tfrac12$\textup{)},
  \[
    \norm{(A - \lambda I)^{-1}}_2
      \le \norm{(A - \lambda I)^{-1}}_1^{\,1-\theta}\,
          \norm{(A - \lambda I)^{-1}}_\infty^{\,\theta},
  \]
  and one may take the infimum of the right-hand side over the two computed
  bounds.
\end{theorem}

\begin{proof}
The Riesz--Thorin theorem gives $\norm{M}_2 \le \norm{M}_1^{1-\theta}
\norm{M}_\infty^{\theta}$ for any matrix $M$ and $\theta \in [0,1]$ with the
stated exponent relation; at $\theta = \tfrac12$ this is the geometric mean.
Apply it to $M = (A - \lambda I)^{-1}$ and insert the bounds of
Theorem~\ref{thm:resolvent}.
\end{proof}

\subsubsection*{Sharper, per-row resolvent bounds}

The bounds above collapse the $\lambda$-dependence into the single scalar
$\min_i\abs{\lambda - \tilde d_i}$, discarding the fact that the $i$-th row of
$M = (\tD - \lambda I)^{-1}S$ is divided by its \emph{own} distance $\abs{\lambda
- \tilde d_i}$.  Retaining the row structure gives strictly sharper bounds.  The
cleanest route observes that the entire $\lambda$-dependence collapses into one
matrix: since $(\tD - \lambda I)M = S$,
\[
  (I - M)^{-1}(\tD - \lambda I)^{-1}
    = \bigl[(\tD - \lambda I)(I - M)\bigr]^{-1}
    = (\tD - \lambda I - S)^{-1},
\]
so that \eqref{eq:resolvent-factored} becomes
\begin{equation}\label{eq:resolvent-clean}
  (A - \lambda I)^{-1} = \tX\,(\tD - S - \lambda I)^{-1}\,\tX^{-1}.
\end{equation}
The middle factor is the resolvent of the explicit, $\lambda$-independent matrix
$\tD - S$, to which one applies Gerschgorin/Varga directly.

\begin{theorem}{Recentered resolvent bound}{resolvent-row}
  Let $S = (I+R_2)^{-1}R_1$ and suppose $\norm{R_2}_\infty < 1$.  Define the
  recentered gaps
  \[
    \hat g_i(\lambda) \coloneqq \abs{\lambda - (\tilde d_i - S_{ii})}
      - \sum_{j \ne i}\abs{S_{ij}},
    \qquad
    \hat g_\infty(\lambda) \coloneqq \min_i \hat g_i(\lambda).
  \]
  If $\hat g_\infty(\lambda) > 0$ then $\lambda \notin \sigma(A)$ and
  \[
    \norm{(A - \lambda I)^{-1}}_\infty
      \;\le\; \frac{\norm{\tX}_\infty\,\norm{\tX^{-1}}_\infty}{\hat g_\infty(\lambda)}
      \;\le\; \frac{\norm{\tX}_\infty\,\norm{Y}_\infty}
                   {(1 - \norm{R_2}_\infty)\,\hat g_\infty(\lambda)}.
  \]
  The $\ell^1$ analogue holds with column sums, $\hat g^{(1)}_j(\lambda) =
  \abs{\lambda - (\tilde d_j - S_{jj})} - \sum_{i \ne j}\abs{S_{ij}}$, and the
  two interpolate to an $\ell^2$ bound exactly as in
  Theorem~\textup{\ref{thm:resolvent2}}.
\end{theorem}

\begin{proof}
By \eqref{eq:resolvent-clean} it suffices to bound $\norm{(\tD - S - \lambda
I)^{-1}}_\infty$.  The matrix $\tD - S - \lambda I$ has diagonal entries
$(\tilde d_i - S_{ii}) - \lambda$ and off-diagonal entries $-S_{ij}$; its
$i$-th Gerschgorin disc is centered at $\lambda - (\tilde d_i - S_{ii})$ (as a
function of the shift) with radius $\sum_{j\ne i}\abs{S_{ij}}$.  When $\hat
g_\infty(\lambda) > 0$ the matrix is strictly diagonally dominant after the
shift, hence nonsingular, with $\norm{(\tD - S - \lambda I)^{-1}}_\infty \le
1/\hat g_\infty(\lambda)$ by the standard Varga bound on the inverse of a
diagonally dominant matrix.  Insert this and $\norm{\tX^{-1}}_\infty \le
\norm{Y}_\infty/(1 - \norm{R_2}_\infty)$ into \eqref{eq:resolvent-clean}.
\end{proof}

\begin{corollary}{Computed form of the recentred bound}{resolvent-row-computed}
  With $\tilde S$, $\beta_\infty$ as in Lemma~\ref{lem:Sresidual}, define the
  computed recentred gap
  \[
    \tilde g_i(\lambda) \coloneqq \abs{\lambda - (\tilde d_i - \tilde S_{ii})}
      - \sum_{j\ne i}\abs{\tilde S_{ij}} - \beta_\infty,
    \qquad
    \tilde g_\infty(\lambda) \coloneqq \min_i \tilde g_i(\lambda).
  \]
  If $\tilde g_\infty(\lambda) > 0$ then $\hat g_\infty(\lambda) \ge \tilde
  g_\infty(\lambda) > 0$, so $\lambda \notin \sigma(A)$ and
  \[
    \norm{(A - \lambda I)^{-1}}_\infty
      \le \frac{\norm{\tX}_\infty\,\norm{Y}_\infty}
               {(1 - \norm{R_2}_\infty)\,\tilde g_\infty(\lambda)}.
  \]
  Every quantity is computed from $\tX, Y, R_1, R_2, \tilde S$; no inverse is
  formed.
\end{corollary}

\begin{proof}
By Lemma~\ref{lem:Sresidual}, $\sum_j\abs{S_{ij} - \tilde S_{ij}} \le
\beta_\infty$ for each $i$, so $\abs{(\tilde d_i - S_{ii}) - (\tilde d_i -
\tilde S_{ii})} + \sum_{j\ne i}\bigl|\abs{S_{ij}} - \abs{\tilde
S_{ij}}\bigr| \le \beta_\infty$.  Hence $\hat g_i(\lambda) \ge \tilde
g_i(\lambda)$ for every $i$, and the claim follows from
Theorem~\ref{thm:resolvent-row}.
\end{proof}

\begin{corollary}{Computed recentred bound, $\ell^1$}{resolvent-row-computed-1}
  With $\tilde S$, $\beta_1 = \norm{(I+R_2)\tilde S - R_1}_1/(1 - \norm{R_2}_1)$,
  define the computed recentred \emph{column} gap
  \[
    \tilde g^{(1)}_j(\lambda) \coloneqq \abs{\lambda - (\tilde d_j - \tilde S_{jj})}
      - \sum_{i\ne j}\abs{\tilde S_{ij}} - \beta_1,
    \qquad
    \tilde g_1(\lambda) \coloneqq \min_j \tilde g^{(1)}_j(\lambda).
  \]
  If $\norm{R_2}_1 < 1$ and $\tilde g_1(\lambda) > 0$, then $\lambda \notin
  \sigma(A)$ and
  \[
    \norm{(A - \lambda I)^{-1}}_1
      \le \frac{\norm{\tX}_1\,\norm{Y}_1}
               {(1 - \norm{R_2}_1)\,\tilde g_1(\lambda)}.
  \]
\end{corollary}

\begin{proof}
Identical to Corollary~\ref{cor:resolvent-row-computed} with rows replaced by
columns: by Lemma~\ref{lem:Sresidual} in the $1\to 1$ norm, $\sum_i\abs{S_{ij}
- \tilde S_{ij}} \le \beta_1$ for each column $j$, so the recentred column gap
of $\tD - S - \lambda I$ degrades by at most $\beta_1$.  Apply the Varga bound
to $(\tD - S - \lambda I)^{-1}$ in the $1\to 1$ norm (column diagonal
dominance) and insert $\norm{\tX^{-1}}_1 \le \norm{Y}_1/(1 - \norm{R_2}_1)$ into
\eqref{eq:resolvent-clean}.
\end{proof}

\begin{corollary}{Computed recentred bound, $\ell^2$ by interpolation}{resolvent-row-computed-2}
  Under the hypotheses of both Corollaries~\textup{\ref{cor:resolvent-row-computed}}
  and~\textup{\ref{cor:resolvent-row-computed-1}} \textup{(}so $\tilde
  g_\infty(\lambda) > 0$ and $\tilde g_1(\lambda) > 0$\textup{)},
  \[
    \norm{(A - \lambda I)^{-1}}_2
      \le \sqrt{\,\frac{\norm{\tX}_1\norm{Y}_1}{(1-\norm{R_2}_1)\,\tilde g_1(\lambda)}
                 \cdot
                 \frac{\norm{\tX}_\infty\norm{Y}_\infty}{(1-\norm{R_2}_\infty)\,\tilde g_\infty(\lambda)}}\,,
  \]
  and more generally
  $\norm{(A-\lambda I)^{-1}}_2 \le \bigl[\text{$\ell^1$ bound}\bigr]^{1-\theta}
  \bigl[\text{$\ell^\infty$ bound}\bigr]^{\theta}$ for $\theta\in[0,1]$, of which
  one takes the sharpest.
\end{corollary}

\begin{proof}
Riesz--Thorin applied to $(A - \lambda I)^{-1}$, inserting the bounds of
Corollaries~\ref{cor:resolvent-row-computed}
and~\ref{cor:resolvent-row-computed-1}.
\end{proof}

\begin{remark}{Three levels of sharpness}{threelevels}
  Writing $\rho_i = \sum_j\abs{S_{ij}}$ for the full row sum, the denominators
  form a chain of increasing sharpness:
  \[
    \underbrace{\min_i\abs{\lambda - \tilde d_i} - \norm{S}_\infty}_{\text{uniform (Thm.~\ref{thm:resolvent})}}
    \;\le\;
    \underbrace{\min_i\bigl(\abs{\lambda - \tilde d_i} - \rho_i\bigr)}_{\text{per-row, original centres}}
    \;\le\;
    \underbrace{\min_i\bigl(\abs{\lambda - (\tilde d_i - S_{ii})} - \!\!\sum_{j\ne i}\!\abs{S_{ij}}\bigr)}_{\text{recentred (Thm.~\ref{thm:resolvent-row})}}.
  \]
  Each is at least the previous, so the corresponding resolvent bounds decrease.
  The recentred form uses the diagonal of $S$ to shift each centre $\tilde d_i
  \to \tilde d_i - S_{ii}$ --- a free first-order eigenvalue correction --- and
  charges only the off-diagonal row sums as radii.  All three share the same
  conditioning prefactor $\kappa_\infty(\tX)$ and are computable from $\tX, Y,
  R_1, R_2$.
\end{remark}

\begin{remark}{Conditioning is unavoidable here, but not in the enclosure}{condrole}
  Theorem~\ref{thm:diag} (the enclosure) carries \emph{no} factor
  $\kappa_p(\tX)$: deciding whether $\lambda$ is an eigenvalue is
  scale-invariant.  The resolvent \emph{norm}, by contrast, must see the
  non-normality of $A$, and $\kappa_p(\tX)$ is exactly the Bauer--Fike
  conditioning that quantifies it.  Its reappearance here is therefore correct
  and necessary, not an artefact --- in contrast to a stray $\sqrt n$, which
  the interpolated bound avoids by using the genuine $\ell^1$ and $\ell^\infty$
  condition numbers rather than a norm-equivalence detour through $\ell^2$.
\end{remark}



\begin{remark}{Why this is cheaper than a transpose-residual bound}{cheap}
  A direct bound built from $\tX^T A \tX - \tD$ and $I - \tX^T\tX$ must carry
  the conditioning constants $\norm{\tX}_\infty$ and $\norm{\tX^T}_\infty$ (the
  latter is \emph{not} $1$ even for orthogonal $\tX$, and can be of order
  $\sqrt n$), together with an explicit $\norm{A}_\infty$ factor.  Miyajima's
  formulation hides all of this inside the single computed residual $R_2 = Y\tX
  - I$: the quality of the approximate inverse $Y \approx \tX^{-1}$ is what
  makes $\norm{R_2}_\infty$ small, and no separate $\norm{A}_\infty$ or
  conditioning factor appears.  Only two matrix products, $R_1$ and $R_2$, are
  formed.
\end{remark}

\begin{remark}{Role of $Y$}{roleY}
  The matrix $Y$ is free.  The bound is valid for \emph{any} $Y$ with
  $\norm{R_2}_\infty < 1$, and is sharpest when $Y$ is close to $\tX^{-1}$.
  Computing $Y$ as a floating-point inverse of $\tX$ and then enclosing $R_1,
  R_2$ rigorously is the standard recipe; nonsingularity of $\tX$ need not be
  assumed beforehand, since $\norm{R_2}_\infty < 1$ verifies it.
\end{remark}

\begin{remark}{Bauer--Fike and Gerschgorin}{bf}
  Theorem~\ref{thm:diag} is the verified Bauer--Fike theorem: in exact
  arithmetic with $Y = \tX^{-1}$ and $\tX$ diagonalizing $A$ exactly, $R_1 = 0$
  and $\varepsilon = 0$.  Corollary~\ref{cor:individual} is its Gerschgorin
  refinement: taking $\tX = I$, $Y = I$, $\tD = \operatorname{diag}(A)$ gives $S
  = A - \tD$ and $\rho_i = \sum_{j\ne i}\abs{A_{ij}}$, the classical discs.
\end{remark}

\begin{remark}{No symmetry used}{symmetry}
  Nothing above uses symmetry or normality of $A$; the discs are complex.
  Symmetry enters only afterwards, to intersect the discs with $\R$.
\end{remark}



\section{Approximate Schur form}

\subsection{Setup}

We now relax the requirement that $\tD$ be diagonal.  Suppose the eigensolver
returns an approximate \emph{Schur} factorization: an approximate invariant
basis $\tX$ and an upper triangular (or, in the real case, quasi-upper
triangular) matrix $\tT$ with $A\tX \approx \tX\tT$.  The diagonal of $\tT$
carries the approximate eigenvalues $\tilde t_1, \ldots, \tilde t_n$.  All other
data are exactly as in Section~1: an auxiliary $Y \approx \tX^{-1}$, the two
residuals
\[
  R_1 \coloneqq Y(A\tX - \tX\tT), \qquad R_2 \coloneqq Y\tX - I,
\]
a computed $\tilde S \approx S \coloneqq (I+R_2)^{-1}R_1$ with residual $T =
(I+R_2)\tilde S - R_1$ and slacks $\beta_p = \norm{T}_p/(1 - \norm{R_2}_p)$.

The algebraic identity of Section~1 is unchanged.  For $\lambda \notin
\sigma(\tT)$,
\begin{equation}\label{eq:schur-collapse}
  (A - \lambda I)^{-1} = \tX\,(\tT - S - \lambda I)^{-1}\,\tX^{-1},
\end{equation}
the only difference being that $\tT - S - \lambda I$ is now quasi-triangular
rather than diagonal-plus-small.

\subsection{Where the difficulty goes}

Running the singularity argument of Theorem~\ref{thm:diag} verbatim gives, at an
eigenvalue $\lambda$ of $A$,
\[
  1 \le \norm{(\tT - \lambda I)^{-1}}_p\,\norm{S}_p,
  \qquad\text{i.e.}\qquad
  \norm{(\tT - \lambda I)^{-1}}_p^{-1} \le \norm{S}_p.
\]
This is a clean statement, but it has merely \emph{relocated} the problem into
bounding the triangular resolvent $\norm{(\tT - \lambda I)^{-1}}_p$.  Whether
that is easy depends entirely on the norm.

\begin{lemma}{Triangular resolvent in $\ell^\infty$ and $\ell^1$}{triag}
  Let $\tT$ be upper triangular with diagonal $\tilde t_1, \ldots, \tilde t_n$.
  For $\lambda$ with
  \[
    g_\infty(\lambda) \coloneqq \min_i\Bigl(\abs{\lambda - \tilde t_i}
      - \sum_{j>i}\abs{\tT_{ij}}\Bigr) > 0,
  \]
  the matrix $\tT - \lambda I$ is strictly diagonally dominant by rows, hence
  $\norm{(\tT - \lambda I)^{-1}}_\infty \le 1/g_\infty(\lambda)$.  Analogously,
  if
  \[
    g_1(\lambda) \coloneqq \min_j\Bigl(\abs{\lambda - \tilde t_j}
      - \sum_{i<j}\abs{\tT_{ij}}\Bigr) > 0,
  \]
  then $\norm{(\tT - \lambda I)^{-1}}_1 \le 1/g_1(\lambda)$.
\end{lemma}

\begin{proof}
Diagonal dominance: the $i$-th row of $\tT - \lambda I$ has diagonal entry
$\tilde t_i - \lambda$ and off-diagonal absolute sum $\sum_{j>i}\abs{\tT_{ij}}$
(only $j > i$ by upper-triangularity); the stated Varga bound on the inverse of
a diagonally dominant matrix gives the $\ell^\infty$ estimate, and the column
version gives the $\ell^1$ estimate.
\end{proof}

So in $\ell^\infty$ and $\ell^1$ the triangular resolvent is bounded in closed
form by a diagonal entry minus an off-diagonal row (resp.\ column) sum --- the
elementary content of Gerschgorin for triangular matrices.  No square roots, no
singular values.

\subsection{The free triangular plot as a first pseudospectral picture}

The gap $g_\infty(z)$ of Lemma~\ref{lem:triag} is, after $\tT$ is computed,
\emph{almost free} to evaluate: the off-diagonal row tails $\tau_i \coloneqq
\sum_{j>i}\abs{\tT_{ij}}$ are fixed once and for all, so at each test point $z$
\[
  g_\infty(z) = \min_i\bigl(\abs{z - \tilde t_i} - \tau_i\bigr)
\]
costs $O(n)$ scalar operations --- no linear algebra, no residual, no inverse.
One may therefore sweep a fine grid in $z$ and plot the level sets of $g_\infty$
directly.  By Lemma~\ref{lem:triag}, $g_\infty(z)$ is a certified lower bound on
$\sigma_{\min}(\tT - zI)$, so
\[
  \bigl\{ z : g_\infty(z) > \varepsilon \bigr\}
  \quad\text{is a rigorous \emph{exclusion region} for }
  \sigma_\varepsilon(\tT),
\]
and its level set $g_\infty(z) = \varepsilon$ is a rigorous \emph{outer}
boundary of the $\varepsilon$-pseudospectrum of $\tT$ (transported to $A$ by the
collapse \eqref{eq:schur-collapse} and the conditioning $\kappa(\tX)$, with the
residual $S$ entering only as the $O(\norm{R_1})$ correction of
Section~3.3).

Three points make this the natural starting layer of the whole computation.

\begin{itemize}
  \item \emph{Rigorous from the first plot.}  The standard pseudospectrum
    picture evaluates $\sigma_{\min}(A - zI)$ approximately on a grid by an SVD
    at each point, with no error control.  The triangular gap instead returns a
    \emph{certified} lower bound on $\sigma_{\min}$ at every grid point, so the
    very first plot already carries guarantees --- loose, but rigorous.
  \item \emph{Almost free.}  The SVD-per-point cost $O(n^3)$ that makes fine
    pseudospectral grids expensive is replaced by $O(n)$ per point.  High
    resolution is immediate.
  \item \emph{It says where to spend.}  Where $g_\infty(z)$ is comfortably
    positive, exclusion is already certified and nothing more is needed.  Where
    it is marginal or negative --- near the spectrum, and wherever the
    off-diagonal mass $\tau_i$ is large --- the cheap bound is inconclusive, and
    \emph{only there} does one escalate to the recentred gap
    (Theorem~\ref{thm:resolvent-row}), then the interpolated $\ell^2$ gap, then
    the cluster $\sigma_{\min}$ assembly (Theorem~\ref{thm:assembly}).  The free
    plot is thus also the refinement trigger of the adaptive strategy of
    Section~2.4.
\end{itemize}

The cost split is then sharp.  \emph{Setup}, once and $O(n^3)$: the basis
$\tX$, the approximate inverse $Y$, the products $B = Y\tX$ (whence $R_2 = B -
I$) and $C = YA\tX$, and the reduced solve $\tilde S$.  \emph{Per test point
$z$}: the free triangular gap $g_\infty(z) = \min_i(\abs{z - \tilde t_i} -
\tau_i)$ is $O(n)$ and uses none of the residual data; the refined exclusion
certificate uses the affine update $R_1(z) = C - zB$ at $O(n^2)$, followed by an
$O(n^2)$ gap or row-sum evaluation.  No matrix product occurs inside the $z$
loop.  A grid of $G$ points therefore costs $O(n^3) + O(G\,n)$ for the free
layer and $O(n^3) + O(G\,n^2)$ for the refined layer --- against $O(G\,n^3)$ for
the uncertified SVD-per-point standard.

\begin{remark}{The free plot is an honest outer bound, loose where it matters}{looseplot}
  Because $g_\infty(z)$ subtracts the full row tail $\tau_i$ from every centre,
  it is pessimistic exactly where the off-diagonal mass is large: near clustered
  or strongly non-normal eigenvalues the gap goes negative well outside the true
  pseudospectrum, so the free plot \emph{over}-estimates $\sigma_\varepsilon$
  there.  This is the correct behaviour for a rigorous outer enclosure --- it
  never misses pseudospectrum --- but the picture must not be read as tight: the
  refinement layers exist precisely to contract this overshoot in the region
  that matters, while the free plot already settles the (typically large)
  benign region at negligible cost.
\end{remark}

\begin{lemma}{Local sandwich: a free, per-point overestimate certificate}{sandwich}
  Let $\tT = D + N$ be upper triangular, $N$ strictly upper triangular, and for
  each $i$ let $\tau_i = \sum_{j>i}\abs{\tT_{ij}}$.  For any $z$, let $i^*(z)$
  attain the minimum in $g_\infty(z) = \min_i(\abs{z - \tilde t_i} - \tau_i)$.
  Then
  \[
    0 \;\le\; \sigma_{\min}(\tT - zI) - g_\infty(z) \;\le\; \tau_{i^*(z)}
      \;\le\; \max_i \tau_i.
  \]
\end{lemma}

\begin{proof}
The lower inequality is Lemma~\ref{lem:triag} ($g_\infty(z) \le
\sigma_{\min}(\tT - zI)$).  For the upper one, $\sigma_{\min}(\tT - zI) \le
\abs{z - \tilde t_{i}}$ for every $i$ (the smallest singular value is at most the
distance to any eigenvalue, here $\tilde t_i$, since $\tT - zI$ is triangular
with that diagonal entry), so in particular $\sigma_{\min}(\tT - zI) \le \abs{z
- \tilde t_{i^*}} = g_\infty(z) + \tau_{i^*}$.  Subtract $g_\infty(z)$.
\end{proof}

The estimate is local, free, and sharp where it matters: the overshoot of the
free gap at $z$ is controlled by the tail $\tau_{i^*(z)}$ of the single active
row, not by any global quantity, and $\tau_{i^*(z)}$ is already known once the
minimising index is found.  (Globally, $\max_i\tau_i$ is a row-wise departure
from normality, comparable to the classical Frobenius measure but without its
dimensional penalty; we do not need it.)  The same number routes refinement: if
$\tau_{i^*(z)} \ll \varepsilon$ the free value is essentially exact at $z$ and no
refinement is warranted; if $\tau_{i^*(z)} \gtrsim \varepsilon$ the free plot is
loose there and one escalates.

\subsection{Lipschitz cell certificates and box bisection}

The map $z \mapsto \sigma_{\min}(A - zI)$ is $1$-Lipschitz, which upgrades a
single certified evaluation into a certificate over a whole region and turns the
sweep from blind grid sampling into a rigorous, contour-concentrated
subdivision.

\begin{lemma}{Lipschitz cell certificate}{lipschitz}
  For all $z, z' \in \C$,
  \[
    \abs{\sigma_{\min}(A - zI) - \sigma_{\min}(A - z'I)} \le \abs{z - z'}.
  \]
  Consequently, if $Z$ is a square cell with centre $z_0$ and half-diagonal $r$
  \textup{(}so $\abs{z - z_0}\le r$ for all $z\in Z$\textup{)}, and
  $\underline{\sigma}(z_0)$ is any certified lower bound on $\sigma_{\min}(A -
  z_0 I)$, then
  \[
    \sigma_{\min}(A - zI) \ge \underline{\sigma}(z_0) - r
    \qquad\text{for every } z \in Z.
  \]
  In particular, if $\underline{\sigma}(z_0) - r > \varepsilon$ then $Z \cap
  \sigma_\varepsilon(A) = \varnothing$: the entire cell is certified outside the
  $\varepsilon$-pseudospectrum from the single centre evaluation.
\end{lemma}

\begin{proof}
By Weyl's perturbation bound for singular values, $\sigma_{\min}$ is
$1$-Lipschitz in the matrix with respect to the spectral norm, and $(A - zI) -
(A - z'I) = (z' - z)I$ has spectral norm $\abs{z - z'}$.  The cell statement
follows with $z' = z_0$ and $\abs{z - z_0}\le r$.
\end{proof}

The constant is exact and dimension-independent --- it is literally $\abs{z -
z_0}$, with no Lipschitz constant to estimate --- so the cell certificate costs
nothing beyond the centre evaluation and the cell geometry.  Symmetrically, an
\emph{upper} bound $\overline{\sigma}(z_0)$ from the inclusion certificate
(Section~3) gives $\sigma_{\min}(A - zI) \le \overline{\sigma}(z_0) + r$ on $Z$,
so $\overline{\sigma}(z_0) + r < \varepsilon$ certifies the whole cell
\emph{inside} $\sigma_\varepsilon(A)$.

\paragraph{Box bisection.}
This yields a quadtree sweep that is a rigorous enclosure of the continuum, not
a sampling of it.  Evaluate the centre of a cell with the cheapest certificate
(the free triangular gap):
\begin{itemize}
  \item if $\underline{\sigma}(z_0) - r > \varepsilon$, mark the cell
    \emph{excluded} and stop;
  \item if $\overline{\sigma}(z_0) + r < \varepsilon$, mark it \emph{included}
    and stop;
  \item otherwise the cell is \emph{ambiguous}: the interval
    $[\underline{\sigma}(z_0) - r,\ \overline{\sigma}(z_0) + r]$ straddles
    $\varepsilon$.  Either escalate the certificate at $z_0$ (recentred,
    interpolated $\ell^2$, or cluster --- guided by the local tail
    $\tau_{i^*(z_0)}$ of Lemma~\ref{lem:sandwich}: a large tail means the free
    value is loose and escalation will help; a small tail means the ambiguity is
    real and the cell must be subdivided), or quadrisect and recurse, until a
    target resolution is reached.
\end{itemize}
Because only ambiguous cells subdivide, refinement concentrates on the
$\varepsilon$-contour --- a one-dimensional set --- so the number of cells scales
with the contour length over the resolution, not the area over the resolution
squared.  The benign region is dispatched by a few large cells, each with one
$O(n)$ free evaluation; the expensive certificates run only in the thin annulus
around the contour.  This is the spatial counterpart of the certificate ladder:
cheap and coarse almost everywhere, expensive and fine only where the geometry
forces it.

\begin{remark}{The $\ell^2$ case is the hard one}{hardell2}
  In $\ell^2$, $\norm{(\tT - \lambda I)^{-1}}_2 = 1/\sigma_{\min}(\tT - \lambda
  I)$, and for a triangular matrix $\sigma_{\min}$ bears \emph{no} relation to
  the diagonal entries: a triangular matrix may have order-one diagonal yet be
  arbitrarily close to singular, the off-diagonal mass conspiring across the
  whole matrix.  There is no diagonal-dominance shortcut.  Certifying
  $\sigma_{\min}(\tT - \lambda I)$ directly is the genuinely hard problem ---
  the pseudospectral / SCI-flavoured difficulty at the heart of non-normal
  spectral computation.  We do not attempt it.  Instead we obtain the $\ell^2$
  bound by interpolation from the two easy norms, accepting the resulting slack
  as the price of avoiding a hard singular-value certificate.
\end{remark}

\begin{remark}{No simplicity or diagonalizability is assumed}{nodefect}
  The entire $\ell^1$/$\ell^\infty$ construction --- Lemma~\ref{lem:triag}, the
  far gap, the localization, and the cheap far block of the assembly
  Theorem~\ref{thm:assembly} --- rests solely on diagonal dominance, a statement
  about the \emph{entries} of $\tT - \lambda I$.  It refers nowhere to the
  distinctness of the $\tilde t_i$, to an eigenvector basis, or to spectral
  projectors.  Consequently every bound here is valid for arbitrary triangular
  $\tT$, including repeated, multiple, and \emph{defective} eigenvalues: a
  Jordan block is strictly diagonally dominant away from its eigenvalue exactly
  as a diagonalizable one is, and the resolvent bound $1/g_p(\lambda)$ correctly
  blows up as $\lambda$ approaches the (possibly defective) diagonal value.  The
  only object sensitive to fine spectral structure is the cluster
  $\sigma_{\min}(\tT_{\mathcal{C}\mathcal{C}} - \lambda I)$ --- and, as the
  nilpotent example $\sigma_{\min}(J - \lambda I) \sim \abs{\lambda}^{m}$ shows,
  $\sigma_{\min}$ is precisely the quantity that reports defectiveness
  correctly, where an eigenvalue-distance would not.  This distinguishes the
  present bounds from perturbation-theoretic ones (Bauer--Fike with an
  eigenvector condition number, or Davis--Kahan), which require diagonalizability
  and degrade or fail on defective spectra.
\end{remark}

\subsection{Enclosure and resolvent via the collapsed matrix}

Rather than bound $\norm{(\tT - \lambda I)^{-1}}$ and $\norm{S}$ separately, it
is sharper to apply Lemma~\ref{lem:triag} directly to the collapsed
quasi-triangular matrix $\tT - S - \lambda I$ of \eqref{eq:schur-collapse},
recentring by the diagonal of $S$ exactly as in
Theorem~\ref{thm:resolvent-row}.  Write $\hat t_i \coloneqq \tilde t_i - S_{ii}$
for the recentred eigenvalue estimates and, in computed form, $\check t_i
\coloneqq \tilde t_i - \tilde S_{ii}$.

\begin{theorem}{Schur resolvent bound, $\ell^\infty$}{schur-resolvent}
  Let $\norm{R_2}_\infty < 1$ and define the recentred quasi-triangular gap
  \[
    G_\infty(\lambda) \coloneqq \min_i\Bigl(\abs{\lambda - \hat t_i}
      - \sum_{j \ne i}\abs{(\tT - S)_{ij}}\Bigr),
  \]
  where the sum runs over the off-diagonal of $\tT - S$ in row $i$.  If
  $G_\infty(\lambda) > 0$ then $\lambda \notin \sigma(A)$ and
  \[
    \norm{(A - \lambda I)^{-1}}_\infty
      \le \frac{\norm{\tX}_\infty\,\norm{Y}_\infty}
               {(1 - \norm{R_2}_\infty)\,G_\infty(\lambda)}.
  \]
  In computed form, replacing $S$ by $\tilde S$ degrades $G_\infty$ by at most
  $\beta_\infty$: with
  \[
    \check G_\infty(\lambda) \coloneqq \min_i\Bigl(\abs{\lambda - \check t_i}
      - \sum_{j\ne i}\abs{(\tT - \tilde S)_{ij}}\Bigr) - \beta_\infty,
  \]
  if $\check G_\infty(\lambda) > 0$ the same bound holds with $G_\infty$
  replaced by $\check G_\infty$.
\end{theorem}

\begin{proof}
By \eqref{eq:schur-collapse} it suffices to bound $\norm{(\tT - S - \lambda
I)^{-1}}_\infty$.  The matrix $\tT - S - \lambda I$ has diagonal entries
$\hat t_i - \lambda$ and off-diagonal row sums $\sum_{j\ne i}\abs{(\tT -
S)_{ij}}$; when $G_\infty(\lambda) > 0$ it is strictly diagonally dominant,
so $\norm{(\tT - S - \lambda I)^{-1}}_\infty \le 1/G_\infty(\lambda)$ by
Lemma~\ref{lem:triag} (applied to the recentred matrix).  Insert this and
$\norm{\tX^{-1}}_\infty \le \norm{Y}_\infty/(1 - \norm{R_2}_\infty)$ into
\eqref{eq:schur-collapse}.  For the computed form, Lemma~\ref{lem:Sresidual}
gives $\sum_j\abs{S_{ij} - \tilde S_{ij}} \le \beta_\infty$ per row, so the
recentred gap degrades by at most $\beta_\infty$.
\end{proof}

The $\ell^1$ version is identical with column sums, $\check G_1(\lambda) =
\min_j(\abs{\lambda - \check t_j} - \sum_{i\ne j}\abs{(\tT - \tilde S)_{ij}}) -
\beta_1$, giving $\norm{(A - \lambda I)^{-1}}_1 \le \norm{\tX}_1\norm{Y}_1 /
[(1 - \norm{R_2}_1)\check G_1(\lambda)]$.

\begin{corollary}{Schur resolvent bound, $\ell^2$ by interpolation}{schur-resolvent-2}
  If $\check G_1(\lambda) > 0$ and $\check G_\infty(\lambda) > 0$, then
  \[
    \norm{(A - \lambda I)^{-1}}_2
      \le \sqrt{\,\frac{\norm{\tX}_1\norm{Y}_1}{(1-\norm{R_2}_1)\,\check G_1(\lambda)}
                 \cdot
                 \frac{\norm{\tX}_\infty\norm{Y}_\infty}{(1-\norm{R_2}_\infty)\,\check G_\infty(\lambda)}}\,.
  \]
\end{corollary}

\begin{proof}
Riesz--Thorin applied to $(A - \lambda I)^{-1}$, inserting the $\ell^1$ and
$\ell^\infty$ bounds.
\end{proof}

\begin{remark}{What the Schur form costs and saves}{schurcost}
  Relative to Section~1, the only change is that the off-diagonal radius in each
  gap now includes the genuine Schur entries $\tT_{ij}$ ($j > i$ in row $i$),
  not merely the residual entries $S_{ij}$.  This off-diagonal mass is the
  non-normality of $A$, made explicit.  In $\ell^1$ and $\ell^\infty$ it is
  harmless --- a finite sum subtracted from each gap.  Only in $\ell^2$ would it
  be dangerous, and there we never confront it directly: the interpolation in
  Corollary~\ref{cor:schur-resolvent-2} produces an $\ell^2$ bound from the two
  triangular estimates, deliberately bypassing the certified smallest singular
  value of Remark~\ref{rem:hardell2}.
\end{remark}



\subsection{An adaptive contour strategy}

The size accounting of the previous remarks suggests a practical division of
labour.  All the expensive objects are computed \emph{once}, up front, and are
independent of the spectral parameter $\lambda$:
\[
  \tX, \quad \tT, \quad Y \approx \tX^{-1}, \quad
  R_1 = Y(A\tX - \tX\tT), \quad R_2 = Y\tX - I, \quad
  \tilde S \approx (I+R_2)^{-1}R_1,
\]
together with the scalar slacks $\beta_\infty, \beta_1$ and the norms
$\norm{\tX}_p, \norm{Y}_p, \norm{R_2}_p$.  This is one factorization, one
approximate inverse, one linear solve, and a handful of matrix products ---
$O(n^3)$ in total, paid once.  Every subsequent operation is $O(n)$ or $O(n^2)$
per value of $\lambda$.

\paragraph{Heuristic pass.}
For each approximate eigenvalue $\tilde t_i$ form the off-diagonal tails
\[
  R_i \coloneqq \sum_{j>i}\abs{\tT_{ij}} \ (\text{row tail}), \qquad
  C_i \coloneqq \sum_{j<i}\abs{\tT_{ji}} \ (\text{column tail}),
\]
and propose the candidate contour
\[
  \Gamma_i \coloneqq \partial B\bigl(\tilde t_i,\ 2\max(R_i, C_i)\bigr).
\]
The factor $2$ is chosen so that, on $\Gamma_i$ and to leading order (the $S$
terms being $O(\norm{R_1})$), both the $\ell^\infty$ and $\ell^1$ gaps satisfy
$\check G_p(\lambda) \gtrsim \max(R_i, C_i)$, keeping the interpolated $\ell^2$
resolvent of order $\kappa_2(\tX)/\max(R_i,C_i)$ --- safely away from
singularity.  Overlaps among the discs $B(\tilde t_i, 2\max(R_i,C_i))$ are read
off in floating point and flag tentative clusters.  This pass is $O(n^2)$ and
not yet rigorous.

\paragraph{Certify or refine.}
For each candidate contour, evaluate the rigorous gaps $\check G_\infty(\lambda),
\check G_1(\lambda)$ of Theorem~\ref{thm:schur-resolvent} along $\Gamma_i$.  Three
outcomes:
\begin{enumerate}
  \item \emph{Gap comfortably positive and $\Gamma_i$ isolated.} The contour is
    certified immediately by Theorem~\ref{thm:schur-resolvent} and
    Corollary~\ref{cor:schur-resolvent-2}; the resolvent is bounded on $\Gamma_i$
    in all of $\ell^1, \ell^2, \ell^\infty$.  Cost $O(n)$ per point.  No
    refinement.
  \item \emph{Discs overlap.} Merge the offending indices into a cluster
    $\mathcal{C}$ and take a single contour around the union, with radius set by
    the collective off-diagonal mass of the cluster.  The block version of the
    bound applies.  Refinement is mere regrouping.
  \item \emph{Gap marginal, or the interpolated $\ell^2$ bound too loose.} Only
    here does one spend: tighten $\tilde S$ by a better solve (shrinking
    $\beta_p$), or escalate to a certified smallest singular value
    $\sigma_{\min}(\tT_{\mathcal{C}} - \lambda I)$ on the small cluster block
    (Remark~\ref{rem:hardell2}).  The hard $\ell^2$ work is thereby confined to
    an $m\times m$ block, $m$ being the cluster multiplicity, never the full $n$.
\end{enumerate}

\begin{remark}{Heuristic and certificate share their quantities}{shared}
  The proposed radius $2\max(R_i, C_i)$ is built from exactly the row and column
  tails $R_i, C_i$ that enter the rigorous gaps $\check G_\infty, \check G_1$.
  Thus the heuristic pass is nothing but the rigorous bound evaluated
  \emph{optimistically} --- without yet checking its positivity hypothesis ---
  and ``refinement'' is checking that hypothesis and, on failure, either
  regrouping (case 2) or escalating to the genuine $\ell^2$ certificate
  (case 3).  No separate heuristic machinery is maintained.  The scheme degrades
  gracefully: a well-separated, nearly normal spectrum is entirely case 1 at
  total cost $O(n^3) + O(n^2)$; non-normality and clustering incur extra cost
  only locally, on the few contours that demand it.
\end{remark}



\section{Block enclosures for clustered spectra}

When approximate eigenvalues cluster, the per-eigenvalue contours of
Section~2.4 overlap and the scalar gaps $\check G_p(\lambda)$ become small or
negative.  The remedy is to treat the cluster as a unit.  This section develops
the block bounds promised in cases 2 and 3 of the adaptive strategy, first for
the complex (triangular) Schur form and then for the real (quasi-triangular)
one.

\subsection{Reordering and partition}

Using \texttt{ordschur} we reorder the Schur form so that the clustered indices
occupy the leading block.  The reordering is an orthogonal (unitary) similarity
and can be certified in ball arithmetic, so the reordered $\tT$ and all the
residual data $R_1, R_2, \tilde S, \beta_p$ remain rigorous.  After reordering,
\[
  \{1, \ldots, n\} = \mathcal{C} \sqcup \mathcal{F},
  \qquad \mathcal{C} = \{1, \ldots, m\}, \quad \mathcal{F} = \{m+1, \ldots, n\},
\]
with $\mathcal{C}$ the cluster and $\mathcal{F}$ the far set.  Because $\tT$ is
upper triangular and $\mathcal{C}$ leads,
\[
  \tT - \lambda I =
  \begin{pmatrix}
    \tT_{\mathcal{C}\mathcal{C}} - \lambda I & \tT_{\mathcal{C}\mathcal{F}} \\
    0 & \tT_{\mathcal{F}\mathcal{F}} - \lambda I
  \end{pmatrix},
\]
the lower-left block vanishing.  We write $w = (w_{\mathcal{C}}, w_{\mathcal{F}})$
for the conformal splitting of a vector.

\subsection{The far gap and localization}

The far block $\tT_{\mathcal{F}\mathcal{F}}$ is itself upper triangular, so its
resolvent is controlled by the elementary triangular bound of
Lemma~\ref{lem:triag}.

\begin{lemma}{Far gap}{fargap}
  Define the $\ell^\infty$ far gap
  \[
    g_{\mathcal{F}}(\lambda) \coloneqq \min_{j\in\mathcal{F}}
      \Bigl(\abs{\lambda - \tilde t_j} - \!\!\sum_{\substack{k>j\\ k\in\mathcal{F}}}\!\!\abs{\tT_{jk}}\Bigr).
  \]
  If $g_{\mathcal{F}}(\lambda) > 0$ then $\tT_{\mathcal{F}\mathcal{F}} - \lambda
  I$ is nonsingular and $\norm{(\tT_{\mathcal{F}\mathcal{F}} - \lambda
  I)^{-1}}_\infty \le 1/g_{\mathcal{F}}(\lambda)$.
\end{lemma}

The cluster bound is meaningful only if the eigenvector mass actually
concentrates on $\mathcal{C}$.  In the present framework this is governed by the
residual $S$ (equivalently its computed proxy $\tilde S$ with slack
$\beta_\infty$), not by a separate orthogonality defect.

\begin{lemma}{Certified localization}{localize}
  Let $u = \tX^{-1}v$ for an eigenvector $Av = \lambda v$, and recall $(\tT -
  \lambda I)u = Su$ from the collapsed identity.  If $g_{\mathcal{F}}(\lambda) >
  0$ and
  \begin{equation}\label{eq:localize-cond}
    \frac{\norm{S}_\infty}{g_{\mathcal{F}}(\lambda)} < 1,
  \end{equation}
  then $\norm{u_{\mathcal{F}}}_\infty < \norm{u}_\infty$, so the dominant
  component lies in the cluster: $\norm{u_{\mathcal{C}}}_\infty =
  \norm{u}_\infty$.  In computed terms \eqref{eq:localize-cond} may be replaced
  by $(\norm{\tilde S}_\infty + \beta_\infty)/g_{\mathcal{F}}(\lambda) < 1$.
\end{lemma}

\begin{proof}
The far rows of $(\tT - \lambda I)u = Su$ read $(\tT_{\mathcal{F}\mathcal{F}} -
\lambda I)u_{\mathcal{F}} = (Su)_{\mathcal{F}}$ (the block $\tT_{\mathcal{F}
\mathcal{C}}$ vanishes).  Hence by Lemma~\ref{lem:fargap},
\[
  \norm{u_{\mathcal{F}}}_\infty
    \le \frac{\norm{(Su)_{\mathcal{F}}}_\infty}{g_{\mathcal{F}}(\lambda)}
    \le \frac{\norm{S}_\infty\,\norm{u}_\infty}{g_{\mathcal{F}}(\lambda)}
    < \norm{u}_\infty
\]
under \eqref{eq:localize-cond}.  Since $\norm{u}_\infty =
\max(\norm{u_{\mathcal{C}}}_\infty, \norm{u_{\mathcal{F}}}_\infty)$, the maximum
is attained on $\mathcal{C}$.  The computed form uses $\norm{S}_\infty \le
\norm{\tilde S}_\infty + \beta_\infty$ (Lemma~\ref{lem:Sresidual}).
\end{proof}

\subsection{Block enclosure, complex Schur}

For bounding the resolvent --- the operative goal in contour-based spectral
computation --- one needs a certified \emph{lower} bound on the cluster block's
smallest singular value, $\sigma_{\min}(\tT_{\mathcal{C}\mathcal{C}} -
S_{\mathcal{C}\mathcal{C}} - \lambda I)$, which is then inverted and assembled
with the cheap far part in Section~3.4.  That lower bound simultaneously
discharges localization: if $\sigma_{\min}(\tT_{\mathcal{C}\mathcal{C}} -
S_{\mathcal{C}\mathcal{C}} - \lambda I) > 0$ then $\lambda \notin
\sigma(\tT_{\mathcal{C}\mathcal{C}} - S_{\mathcal{C}\mathcal{C}})$, so no separate
dominant-component hypothesis is required.  We therefore make the resolvent
assembly of Section~3.4 the central result and record the corresponding
\emph{enclosure} statement --- an \emph{upper} bound on the same
$\sigma_{\min}$, locating where eigenvalues lie --- only as a remark, for use
when one wants to count or place eigenvalues rather than bound the resolvent.

\begin{remark}{Cluster enclosure, when eigenvalue location is wanted}{cluster}
  Suppose $\norm{R_2}_\infty < 1$, $g_{\mathcal{F}}(\lambda) > 0$, and the
  localization condition \eqref{eq:localize-cond} holds.  Then every eigenvalue
  $\lambda$ of $A$ with dominant component in $\mathcal{C}$ satisfies
  \[
    \sigma_{\min}\bigl(\tT_{\mathcal{C}\mathcal{C}} - \lambda I\bigr)
      \le \sqrt{m}\,\Bigl(\norm{S_{\mathcal{C}\mathcal{C}}}_\infty
        + \norm{\tT_{\mathcal{C}\mathcal{F}}}_\infty\,
          \frac{\norm{S}_\infty}{g_{\mathcal{F}}(\lambda)}\Bigr),
  \]
  i.e.\ $\lambda$ lies within that distance of $\sigma(\tT_{\mathcal{C}
  \mathcal{C}})$.  The proof selects the dominant component (in $\mathcal{C}$ by
  Lemma~\ref{lem:localize}), uses the cluster rows $(\tT_{\mathcal{C}\mathcal{C}}
  - \lambda I)u_{\mathcal{C}} = (Su)_{\mathcal{C}} - \tT_{\mathcal{C}\mathcal{F}}
  u_{\mathcal{F}}$ with $\norm{u_{\mathcal{F}}}_\infty \le \norm{S}_\infty
  \norm{u}_\infty / g_{\mathcal{F}}$, then converts the $\ell^\infty$ row
  estimate to $\sigma_{\min}$ at the cost of $\sqrt m$ ($\norm{x}_2 \le \sqrt m
  \norm{x}_\infty$, $\norm{Cx}_2 \ge \sigma_{\min}(C)\norm{x}_2$).  This is an
  \emph{upper} bound on $\sigma_{\min}$: it tells where eigenvalues are, not how
  large the resolvent is, and carries the $\sqrt m$ factor and the
  dominant-component hypothesis --- neither of which the resolvent route of
  Section~3.4 needs.
\end{remark}

\begin{remark}{The cluster $\sigma_{\min}$ is the only hard object, and it is small}{clustersigma}
  The certificate reduces to $\sigma_{\min}(\tT_{\mathcal{C}\mathcal{C}}
  - \lambda I)$ on an $m\times m$ matrix, $m$ the cluster size.  This is the lone
  appearance of the genuinely hard $\ell^2$ singular-value problem of
  Remark~\ref{rem:hardell2}, now confined to a block whose dimension is the
  cluster multiplicity rather than $n$.  For $m \le 2$ it is available in closed
  form (next subsection); for larger $m$ one uses a verified SVD on the small
  block.  The right-hand side is $O(\norm{S}_\infty) = O(\norm{R_1})$ plus a
  coupling term suppressed by $g_{\mathcal{F}}^{-1}$, so a good factorization
  makes the enclosure tight.
\end{remark}

\begin{remark}{Tighter form on $\tT - S$}{tightenTS}
  The theorems above keep $\tT$ and carry $S$ (or $\tilde S$ with slack
  $\beta$) explicitly, which is cleanest to read.  A tighter variant applies the
  whole partition to the collapsed matrix $\tT - S - \lambda I$ of
  \eqref{eq:schur-collapse} instead of $\tT - \lambda I$: the far gap becomes
  $g_{\mathcal{F}}$ computed from $\tT - S$ (recentred by $S_{jj}$ and with the
  $S$ off-diagonal added to the tails), and the residual coupling is folded in
  rather than bounded separately.  All statements hold verbatim with $\tT
  \rightsquigarrow \tT - S$ and $S_{\mathcal{C}\mathcal{C}} \rightsquigarrow 0$
  in the leading term, at the cost of the less transparent recentring.
\end{remark}

\subsection{From the cluster block to the full resolvent}

This is the central result of the section: a certified bound on the full
resolvent $\norm{(A - \lambda I)^{-1}}_2$ requiring only a \emph{lower} bound on
the cluster block's smallest singular value, with no eigenvector selection and
no $\sqrt m$ factor.  The cluster bound controls only the $m\times m$ block
$\tT_{\mathcal{C}\mathcal{C}}$; to reach the full resolvent we assemble it with
the far part through the block-triangular inverse.  Recall the collapse
\eqref{eq:schur-collapse},
\[
  (A - \lambda I)^{-1} = \tX\,(\tT - S - \lambda I)^{-1}\,\tX^{-1},
  \qquad\text{so}\qquad
  \norm{(A - \lambda I)^{-1}}_2 \le \kappa_2(\tX)\,\norm{(\tT - S - \lambda I)^{-1}}_2.
\]
After the cluster reordering, partition the (upper block-triangular)
quasi-triangular matrix
\[
  \tT - S - \lambda I =
  \begin{pmatrix} P & Q \\ 0 & W \end{pmatrix},
  \quad
  P \coloneqq \tT_{\mathcal{C}\mathcal{C}} - S_{\mathcal{C}\mathcal{C}} - \lambda I,
  \quad
  W \coloneqq \tT_{\mathcal{F}\mathcal{F}} - S_{\mathcal{F}\mathcal{F}} - \lambda I,
  \quad
  Q \coloneqq \tT_{\mathcal{C}\mathcal{F}} - S_{\mathcal{C}\mathcal{F}},
\]
the lower-left block being $-S_{\mathcal{F}\mathcal{C}}$, of size $O(\norm{R_1})$,
which we take to vanish to leading order (Remark~\ref{rem:schurcomp} below
carries it).

\begin{theorem}{Full resolvent by block-triangular assembly}{assembly}
  With $P, Q, W$ as above and the lower-left block neglected, suppose $P$ and
  $W$ are nonsingular.  Then
  \[
    \norm{(\tT - S - \lambda I)^{-1}}_2
      \le \bigl(1 + \norm{Q}_2\,\norm{W^{-1}}_2\bigr)\,\norm{P^{-1}}_2
        + \norm{W^{-1}}_2,
  \]
  and hence
  \[
    \norm{(A - \lambda I)^{-1}}_2
      \le \kappa_2(\tX)\Bigl[\bigl(1 + \norm{Q}_2\,\norm{W^{-1}}_2\bigr)\,
        \norm{P^{-1}}_2 + \norm{W^{-1}}_2\Bigr],
  \]
  where the three ingredients are obtained by \emph{different} means matched to
  their cost:
  \begin{itemize}
    \item $\norm{P^{-1}}_2 = 1/\sigma_{\min}(\tT_{\mathcal{C}\mathcal{C}} -
      S_{\mathcal{C}\mathcal{C}} - \lambda I)$ --- the \emph{expensive} certified
      smallest singular value, on the $m\times m$ cluster block only
      \textup{(}closed form for $m\le 2$, verified SVD otherwise\textup{)};
    \item $\norm{W^{-1}}_2$ --- the \emph{cheap} far resolvent, from the
      $\ell^1$/$\ell^\infty$ triangular diagonal-dominance gaps of
      Lemma~\ref{lem:triag} interpolated by Riesz--Thorin,
      $\norm{W^{-1}}_2 \le (g^W_1(\lambda)\,g^W_\infty(\lambda))^{-1/2}$;
    \item $\norm{Q}_2 = \norm{\tT_{\mathcal{C}\mathcal{F}} -
      S_{\mathcal{C}\mathcal{F}}}_2$ --- a single computed block norm.
  \end{itemize}
\end{theorem}

\begin{proof}
For a block upper-triangular matrix with nonsingular diagonal blocks,
\[
  \begin{pmatrix} P & Q \\ 0 & W \end{pmatrix}^{-1}
    = \begin{pmatrix} P^{-1} & -P^{-1}QW^{-1} \\ 0 & W^{-1} \end{pmatrix}.
\]
Bounding each block in the spectral norm and using $\norm{\,\cdot\,}_2 \le
\sum(\text{block norms})$ for the $2\times 2$ block structure gives
$\norm{(\tT - S - \lambda I)^{-1}}_2 \le \norm{P^{-1}}_2 + \norm{P^{-1}}_2
\norm{Q}_2\norm{W^{-1}}_2 + \norm{W^{-1}}_2$, which is the stated bound.  Dress
with $\kappa_2(\tX)$ via the collapse identity.
\end{proof}

The cost asymmetry is the whole point: the far block, of dimension $n - m$, is
handled by the \emph{cheap} triangular gaps; the only \emph{expensive} object
--- a certified $\sigma_{\min}$ --- lives on the $m\times m$ cluster, never on
the full $n\times n$ matrix.  When the cluster is well separated,
$\norm{W^{-1}}_2 = O(1/g_{\mathcal{F}})$ is small, the cross term is negligible,
and the bound collapses to $\norm{(A-\lambda I)^{-1}}_2 \lesssim
\kappa_2(\tX)\,\norm{P^{-1}}_2$ --- the full resolvent is, up to conditioning,
the cluster resolvent.

\begin{remark}{Relation to the Rump--Ogita Schur complement}{schurcomp}
  If the lower-left block $-S_{\mathcal{F}\mathcal{C}}$ is retained \textup{(}or
  if one does not reorder to block-triangular form\textup{)}, the inverse of the
  leading block is replaced by the inverse of the \emph{Schur complement}
  \[
    P_{\mathrm{SC}} = P - Q\,W^{-1}(-S_{\mathcal{F}\mathcal{C}}),
  \]
  and $\norm{P^{-1}}_2$ in Theorem~\ref{thm:assembly} becomes
  $\norm{P_{\mathrm{SC}}^{-1}}_2$.  This is the validated Schur-complement bound
  of Rump--Ogita in full: the conditioning is concentrated in the small block
  $P_{\mathrm{SC}}$, the far block enters only through $W^{-1}$ and the
  coupling.  In our setting the correction $Q W^{-1} S_{\mathcal{F}\mathcal{C}}$
  is quadratic in the residual, of size $O(\norm{R_1}^2/g_{\mathcal{F}})$, hence
  negligible to leading order; this is why the block-triangular form
  \textup{(}with the lower-left block dropped\textup{)} suffices in practice, and
  the genuine quadratic Schur-complement correction is needed only when
  $\norm{R_1}$ is not small.
\end{remark}

\subsection{The cluster block in $\ell^1$ and $\ell^2$}

The $\ell^1$ cluster enclosure is the transpose statement: the far gap uses
column tails, the coupling uses $\norm{\tT_{\mathcal{C}\mathcal{F}}}_1$, the
localization uses $\norm{S}_1$, and the prefactor $\sqrt m$ is unchanged.  The
two interpolate.

\begin{corollary}{Cluster resolvent, $\ell^2$}{cluster2}
  Off the cluster spectrum, with $r_{\mathcal{C}}^{\infty}(\lambda)$ and
  $r_{\mathcal{C}}^{1}(\lambda)$ the $\ell^\infty$ and $\ell^1$ cluster radii,
  the cluster-block resolvent obeys
  \[
    \norm{(\tT_{\mathcal{C}\mathcal{C}} - \lambda I)^{-1}}_2
      \le \frac{1}{\sigma_{\min}(\tT_{\mathcal{C}\mathcal{C}} - \lambda I)},
    \qquad
    \sigma_{\min} \ge \max\Bigl(0,\ \sigma_{\min}^{\mathrm{cert}}\Bigr),
  \]
  where $\sigma_{\min}^{\mathrm{cert}}$ is a verified lower bound on the cluster
  block's smallest singular value, computed directly on the $m\times m$ matrix.
  The full resolvent is recovered through \eqref{eq:schur-collapse} with the
  conditioning prefactor $\kappa_2(\tX) \le \kappa_1(\tX)^{1/2}
  \kappa_\infty(\tX)^{1/2}$.
\end{corollary}

\subsection{The real Schur form}

For real $A$ one uses the real Schur form, in which $\tT$ is
\emph{quasi}-upper-triangular: block upper triangular with diagonal blocks of
size $1$ (real eigenvalues) or $2$ (conjugate complex pairs).  The development
above goes through with ``index'' replaced by ``diagonal block'' throughout.

Partition $\{1,\ldots,n\}$ into the quasi-diagonal blocks $B_1, \ldots, B_p$,
$\abs{B_a}\in\{1,2\}$, with $T_a \coloneqq \tT_{B_a B_a}$ the $a$-th diagonal
block.  The scalar quantity attached to a block is
\[
  s_a(\lambda) \coloneqq \sigma_{\min}(T_a - \lambda I)
    = \mathrm{dist}\bigl(\lambda, \sigma(T_a)\bigr),
\]
exact because $\sigma(T_a)$ is known in closed form (the eigenvalues of a
$2\times 2$ block are the roots of its characteristic quadratic).

\begin{theorem}{Real Schur block enclosure, $\ell^\infty$}{realblock}
  Let $\tT$ be a real Schur form with diagonal blocks $T_1, \ldots, T_p$, and
  let $S = (I+R_2)^{-1}R_1$.  For each block index $a$ define the block-row
  radius
  \[
    \mathcal{R}_a \coloneqq \sum_{b}\norm{(S)_{B_a B_b}}_\infty
      + \sum_{b>a}\norm{\tT_{B_a B_b}}_\infty,
  \]
  the second sum running over the strict block-upper-triangle.  Then every
  eigenvalue $\lambda$ of $A$ satisfies $s_{a^*}(\lambda) \le c_{a^*}\,
  \mathcal{R}_{a^*}$ for some block $a^*$, with $c_a = 1$ if $\abs{B_a}=1$ and
  $c_a = \sqrt 2$ if $\abs{B_a}=2$.  Hence
  \[
    \sigma(A) \subseteq \bigcup_{a=1}^p
      \bigl\{ z \in \C : \sigma_{\min}(T_a - z I) \le c_a \mathcal{R}_a \bigr\}.
  \]
\end{theorem}

\begin{proof}
With $u = \tX^{-1}v$ and the collapsed identity $(\tT - \lambda I)u = Su$,
restrict to the rows of block $B_a$ and use $\tT_{B_a B_b} = 0$ for $b < a$:
\[
  (T_a - \lambda I)u_{B_a}
    = (Su)_{B_a} - \sum_{b > a}\tT_{B_a B_b}u_{B_b}.
\]
The residual term $(Su)_{B_a}$ and the genuine Schur coupling $\sum_{b>a}\tT_{B_a
B_b}u_{B_b}$ are both bounded by the block-row sums times $\norm{u}_\infty$.
Choosing $a^*$ with
$\norm{u_{B_{a^*}}}_\infty = \norm{u}_\infty$ and dividing,
\[
  \norm{(T_{a^*} - \lambda I)u_{B_{a^*}}}_\infty \le \mathcal{R}_{a^*}\norm{u}_\infty.
\]
For the lower bound, $\norm{x}_\infty \ge \tfrac{1}{\sqrt{\abs{B_{a^*}}}}
\norm{x}_2$ on the block and $\norm{Cx}_2 \ge \sigma_{\min}(C)\norm{x}_2$ give
$\norm{(T_{a^*}-\lambda I)u_{B_{a^*}}}_\infty \ge
\tfrac{1}{\sqrt{\abs{B_{a^*}}}}\, s_{a^*}(\lambda)\norm{u_{B_{a^*}}}_2 \ge
\tfrac{1}{\sqrt{\abs{B_{a^*}}}}\,s_{a^*}(\lambda)\norm{u}_\infty$.  Combining and
noting $\abs{B_{a^*}}\le 2$ yields $s_{a^*}(\lambda) \le c_{a^*}\mathcal{R}_{a^*}$.
\end{proof}

\begin{remark}{Explicit $\sigma_{\min}$ for a $2\times 2$ block}{twobytwo}
  For $T_a - \lambda I = \begin{psmallmatrix} a & b \\ c & d \end{psmallmatrix}$
  the singular values are available in closed form.  The cancellation-free
  (Kahan) form is best for validated arithmetic: with
  \[
    p = \sqrt{(a+d)^2 + (b-c)^2}, \qquad q = \sqrt{(a-d)^2 + (b+c)^2},
  \]
  one has $\sigma_{\max} = \tfrac12(p+q)$ and $\sigma_{\min} =
  \tfrac12\abs{p-q}$.  Each of $p, q$ is a square root of a sum of squares, so
  only the outer subtraction $p - q$ requires care, and interval arithmetic
  encloses it tightly even as $\sigma_{\min}\to 0$.  In the shifted block only
  $a + d = (\alpha + \delta) - 2\lambda$ depends on $\lambda$ (for real
  $\lambda$), so $s_a(\lambda)$ is an explicit, cheaply swept function along the
  real axis.  For complex $\lambda$ one takes $\sigma_{\min}^2 =
  \lambda_{\min}\bigl((T_a - \lambda I)^*(T_a - \lambda I)\bigr)$, again a
  $2\times 2$ closed form.
\end{remark}

\begin{remark}{Real clusters and conjugate pairs}{realcluster}
  A conjugate pair is already a $2\times 2$ block --- the smallest possible
  cluster --- handled exactly by Remark~\ref{rem:twobytwo} with no reordering.
  Larger real clusters (several diagonal blocks with nearby spectra) are grouped
  by \texttt{ordschur} into a leading quasi-triangular cluster block
  $\tT_{\mathcal{C}\mathcal{C}}$, and the assembly of
  Theorem~\ref{thm:assembly} applies with $\sigma_{\min}$ certified on that
  small block.  The far gap, localization, and
  $\ell^1/\ell^2$ interpolation are identical to the complex case, with diagonal
  blocks in place of diagonal entries.
\end{remark}

\section{An SVD-frame certificate for $\sigma_{\min}$}
\label{sec:svdframe}

The certificates so far bound $\sigma_{\min}(A - zI)$ through a \emph{surrogate}:
the triangular diagonal-dominance gap of Lemma~\ref{lem:triag} controls
$\norm{(\tT - zI)^{-1}}$ and overshoots by the full off-diagonal mass of the
Schur factor; the recentred and interpolated bounds reach $\ell^2$ only through
the conditioning $\kappa(\tX)$.  We close with a different certificate that
targets $\sigma_{\min}(A - zI)^2$ \emph{directly}, from a single verified
singular value decomposition of $A$, reused across the whole plane.  It carries
no $\kappa$ factor and its overshoot is governed by the misalignment of the left
and right singular vectors rather than by triangular tails --- a genuinely
distinct, and empirically comparable, layer (Remark~\ref{rem:svdvstri}).  It is
also the natural device for the $\sigma_{\min}$ of the small cluster blocks of
Theorem~\ref{thm:assembly} (Remark~\ref{rem:svdcluster}).

\subsection{The normal-equations frame}

Write the Hermitian normal-equations matrix
\[
  M(z) \coloneqq (A - zI)^H(A - zI),
  \qquad \sigma_{\min}(A - zI)^2 = \lambda_{\min}\bigl(M(z)\bigr).
\]
Expanding,
\begin{equation}\label{eq:svdMz}
  M(z) = A^H A - \bar z\,A - z\,A^H + \abs{z}^2 I,
\end{equation}
a fixed quadratic-in-$z$ combination of $A^HA$, $A$, $A^H$ and $I$.  Let $A =
U\Sigma V^H$ be an exact SVD, $U,V$ unitary, $\Sigma =
\operatorname{diag}(\sigma_1,\ldots,\sigma_n)$.  Conjugating \eqref{eq:svdMz} by
the right singular vectors $V$ diagonalizes the constant term, $V^H A^H A V =
\Sigma^2$, and leaves the identity invariant:
\begin{equation}\label{eq:svdideal}
  V^H M(z) V = \bigl(\Sigma^2 + \abs{z}^2 I\bigr)
    - \bigl(\bar z\,W\Sigma + z\,\Sigma W^H\bigr),
  \qquad W \coloneqq V^H U \ \text{(unitary)}.
\end{equation}
The constant and quadratic parts are diagonal --- the $i$-th diagonal entry is
$\sigma_i^2 + \abs{z}^2$, and the $\abs{z}^2$ growth strengthens diagonal
dominance away from the spectrum, so the far region is the easy region.  The
entire off-diagonal, $z$-dependent part is $-(\bar z\,W\Sigma + z\,\Sigma W^H)$,
linear in $z$ with \emph{fixed} coefficient matrices; entrywise, for $i \ne j$,
\[
  [V^H M(z) V]_{ij} = -\bigl(\bar z\,\sigma_j W_{ij}
    + z\,\sigma_i\overline{W_{ji}}\bigr).
\]

\begin{remark}{Where non-normality lives}{svdnonnormal}
  $W = V^H U$ is the misalignment of the right and left singular vectors of $A$.
  $W = I$ iff $U = V$ iff $A$ is Hermitian positive semidefinite, in which case
  \eqref{eq:svdideal} is exactly diagonal and $\lambda_{\min}(M(z))$ is read off
  with no overshoot.  For non-normal $A$ the off-diagonal mass is governed by the
  departure of $W$ from the identity --- an SVD-flavoured measure of
  non-normality that enters only the \emph{linear}-in-$z$ terms, the heavy
  constant part $A^HA$ having been diagonalized exactly.  This is the advantage
  of the SVD frame over a diagonal balancing of $A^HA$, which could render it
  only approximately diagonally dominant.
\end{remark}

\subsection{Certification by a verified residual}

In practice one computes an approximate SVD $A \approx \tU\tSig\tV^H$ in floating
point.  As with $R_1, R_2$ in Section~\ref{sec:diag}, we do not certify the
singular vectors directly but bound, \emph{once}, the residuals
\begin{equation}\label{eq:svdresiduals}
  P \coloneqq \tV^H\tV - I, \qquad
  Q \coloneqq \tU^H\tU - I, \qquad
  E \coloneqq \tU^H A \tV - \tSig,
\end{equation}
the orthogonality defects of the computed factors and the factorization
residual, with certified $2$-norms obtained in interval or ball arithmetic.  All
three are $z$-independent.

\begin{lemma}{SVD-frame form with certified deviation}{svdform}
  Let $\tW \coloneqq \tV^H\tU$.  In the computed right-singular basis,
  \[
    \tV^H M(z) \tV
      = \bigl(\tSig^2 + \abs{z}^2 I\bigr)
        - \bigl(\bar z\,\tW\tSig + z\,\tSig\tW^H\bigr)
        + \mathcal{E}(z),
  \]
  where the deviation from the ideal form \eqref{eq:svdideal} satisfies, with
  $z$-independent constants built from \eqref{eq:svdresiduals},
  \begin{equation}\label{eq:svdEbound}
    \norm{\mathcal{E}(z)}_2
      \;\le\; \underbrace{2\norm{\tSig}_2\norm{E}_2 + \norm{E}_2^2
              + \norm{\tSig}_2^2\norm{Q}_2}_{\text{constant part}}
      \;+\; \abs{z}\,\underbrace{\bigl(2\norm{E}_2 + 2\norm{\tSig}_2\norm{Q}_2\bigr)}_{\text{linear part}}
      \;+\; \abs{z}^2\,\norm{P}_2.
  \end{equation}
\end{lemma}

\begin{proof}[Proof sketch]
Insert $A = \tU\tSig\tV^H + (A - \tU\tSig\tV^H)$ into each term of
\eqref{eq:svdMz} conjugated by $\tV$.  The constant term $\tV^H A^H A\tV =
(\tU^H A\tV)^H(\tU^H A\tV)$ up to the defect $Q$; since $\tU^H A\tV = \tSig + E$,
it equals $\tSig^2 + (\tSig E + E^H\tSig + E^H E)$ modulo
$\norm{\tSig}_2^2\norm{Q}_2$, giving the constant part of \eqref{eq:svdEbound}.
The cross terms $\bar z\,\tV^H A\tV$ and $z\,\tV^H A^H\tV$ equal $\bar
z\,\tW\tSig + z\,\tSig\tW^H$ up to $O(\abs z\,\norm{E}_2)$ and $O(\abs
z\,\norm{\tSig}_2\norm{Q}_2)$, giving the linear part.  The term $\abs z^2\tV^H
\tV = \abs z^2(I + P)$ contributes $\abs z^2\norm{P}_2$.
\end{proof}

\begin{theorem}{Certified $\sigma_{\min}$ surface from one verified SVD}{svdmain}
  With the notation above, for every $z \in \C$,
  \[
    \sigma_{\min}(A - zI)^2 = \lambda_{\min}(M(z))
      \;\ge\; \gamma(z) \coloneqq G(z) - \norm{\mathcal{E}(z)}_2,
  \]
  where $\norm{\mathcal{E}(z)}_2$ is bounded by \eqref{eq:svdEbound} and $G(z)$ is
  the Hermitian Gershgorin lower bound of the ideal form,
  \begin{equation}\label{eq:svdG}
    G(z) \coloneqq \min_{1\le i\le n}\Bigl(
      \underbrace{\tSig_i^2 + \abs{z}^2
        - 2\,\tSig_i\,\Re\!\bigl(\bar z\,\tW_{ii}\bigr)}_{\text{$i$-th diagonal entry}}
      \;-\; \sum_{j\ne i}\bigl|\bar z\,\tSig_j \tW_{ij}
        + z\,\tSig_i\overline{\tW_{ji}}\bigr|
      \Bigr).
  \end{equation}
  Wherever $\gamma(z) > 0$, $z \notin \sigma(A)$ and $\norm{(A - zI)^{-1}}_2 \le
  \gamma(z)^{-1/2}$; wherever $\gamma(z) > \varepsilon^2$, $z$ is certified
  outside the $\varepsilon$-pseudospectrum of $A$.
\end{theorem}

\begin{proof}
$\tV^H M(z)\tV$ is Hermitian and equals $M(z)$ up to a unitary change of basis
modulo the defect $P$, both absorbed in $\mathcal{E}(z)$, so
$\lambda_{\min}(M(z)) \ge \lambda_{\min}(\text{ideal form}) -
\norm{\mathcal{E}(z)}_2$.  For a Hermitian matrix the smallest eigenvalue is
bounded below by the minimal Gershgorin left-endpoint (diagonal entry minus
off-diagonal absolute row sum).  The diagonal entries of the ideal form
\eqref{eq:svdideal} are $\tSig_i^2 + \abs{z}^2 - (\bar z\,\tSig_i\tW_{ii} +
z\,\tSig_i\overline{\tW_{ii}}) = \tSig_i^2 + \abs{z}^2 -
2\tSig_i\Re(\bar z\,\tW_{ii})$ --- the diagonal of the $z$-linear term must be
retained, as it carries the entire bound in the Hermitian-PSD case $\tW = I$ ---
and the off-diagonal entries are $-(\bar z\,\tSig_j\tW_{ij} +
z\,\tSig_i\overline{\tW_{ji}})$, giving \eqref{eq:svdG}.  Subtract
$\norm{\mathcal{E}(z)}_2$.
\end{proof}

\subsection{Cost, reuse, and comparison}

\emph{Once, $O(n^3)$:} the (verified) SVD $\tU,\tSig,\tV$; the residual norms
\eqref{eq:svdresiduals}; the fixed $\tW = \tV^H\tU$ and the coefficient arrays
$\tSig_j\tW_{ij}$, $\tSig_i\overline{\tW_{ji}}$.  \emph{Per $z$, $O(n^2)$:} the
off-diagonal absolute row sums of the $z$-linear part for $G(z)$ and the scalar
polynomial \eqref{eq:svdEbound} for $\norm{\mathcal{E}(z)}_2$; with per-row
envelopes of the linear coefficients this drops to $O(n)$.  No matrix product or
factorization occurs in the $z$ loop --- the same reuse principle as the affine
$R_1(z) = C - zB$ of the bridge theorem.

\begin{remark}{Two independent knobs}{svdknobs}
  The bound $\gamma(z)$ loosens for two separately controlled reasons.
  \emph{Non-normality} enters through $\tW = \tV^H\tU$ in $G(z)$ and vanishes as
  $A \to$ Hermitian-PSD; it is intrinsic to $A$ and cannot be removed by
  recomputation.  \emph{Factorization inaccuracy} enters through
  $\norm{P}_2,\norm{Q}_2,\norm{E}_2$ in $\mathcal{E}(z)$ and vanishes as the SVD
  is computed more accurately (e.g.\ in higher precision); it is a matter of
  effort, not of $A$.  One may therefore tighten the certificate by a more
  accurate SVD without touching the non-normality overshoot, which is present
  even with an exact factorization.
\end{remark}

\begin{remark}{Versus the triangular (Schur) certificate}{svdvstri}
  The free triangular gap of Lemma~\ref{lem:triag}, $\min_i(\abs{z - \tilde t_i}
  - \sum_{j>i}\abs{\tT_{ij}})$, is $O(n)$ and uses only a Schur factorization,
  but it is a surrogate overshooting by the full off-diagonal mass of the Schur
  factor (Lemma~\ref{lem:sandwich}).  The SVD frame targets $\sigma_{\min}^2$
  directly, with overshoot governed by $\tW = \tV^H\tU$ rather than the
  triangular tails, and its far-region behaviour is aided by the $\abs{z}^2$
  diagonal growth.  Which is tighter on a given matrix --- the triangular tails
  $\sum_{j>i}\abs{\tT_{ij}}$ versus the singular-vector misalignment $\tW - I$
  --- is empirical; the point is that the two are not competitors but
  \emph{complementary floors}, combined by the maximum of
  Section~\ref{subsec:svdcombine}.  The SVD frame costs one SVD instead of one
  Schur factorization, comparable.
\end{remark}

\begin{remark}{Cluster blocks}{svdcluster}
  For the $m\times m$ cluster blocks of the block-assembled resolvent bound
  (Theorem~\ref{thm:assembly}), the SVD-frame certificate is a natural choice for
  the required $\sigma_{\min}$: form the verified SVD of the $m\times m$ block
  once and reuse it across the $z$ of that cluster's contour, rather than a fresh
  verified SVD per $z$.
\end{remark}

\begin{remark}{Refinements}{svdtodo}
  Two sharpenings are worth noting.  The constants in \eqref{eq:svdEbound} are
  written generously; keeping $P, Q, E$ as interval/ball matrices rather than
  collapsing to norms tightens the deviation, especially the $\abs{z}^2\norm{P}_2$
  term that dominates far out and competes with the helpful $\abs{z}^2$ diagonal
  growth.  And for the box-bisection sweep one bounds
  $\gamma(z)$ from below over a whole cell by interval evaluation of $G$ and
  $\norm{\mathcal{E}}_2$, exactly as the Lipschitz cell certificate
  (Lemma~\ref{lem:lipschitz}) does for the scalar gaps.
\end{remark}

\subsection{Combining certificates: the maximum of floors}
\label{subsec:svdcombine}

The triangular gap of Lemma~\ref{lem:triag}, the SVD-frame bound
$\gamma(z)^{1/2}$ of Theorem~\ref{thm:svdmain}, and the recentred and cluster
bounds of the earlier sections are all \emph{certified lower bounds} on the same
scalar $\sigma_{\min}(A - zI)$.  Lower bounds compose trivially and rigorously.

\begin{lemma}{Maximum of floors}{maxfloor}
  If $\ell_1(z), \ldots, \ell_k(z)$ are each certified lower bounds on
  $\sigma_{\min}(A - zI)$, then so is $\ell(z) \coloneqq \max_t \ell_t(z)$.
  Consequently $z$ is certified outside $\sigma_\varepsilon(A)$ as soon as
  \emph{any one} $\ell_t(z) > \varepsilon$.
\end{lemma}

There is no hierarchy to respect: the certificates form a family of independent
floors, and one always keeps the largest.  Two of them are cheap and obtained
from a single factorization each.

\paragraph{Schur and SVD, both at setup.}
A Schur factorization and an SVD of $A$ are each a single $O(n^3)$ computation,
negligible against a sweep over thousands of cells.  One therefore computes
\emph{both} once and, at each $z$, takes
\begin{equation}\label{eq:cheapfloor}
  \sigma_{\min}(A - zI) \;\ge\;
    \max\bigl(g_{\mathrm{tri}}(z),\ \gamma(z)^{1/2}\bigr),
\end{equation}
where $g_{\mathrm{tri}}$ is the triangular gap from the Schur factor and
$\gamma^{1/2}$ the SVD-frame bound.  Evaluating both is still $O(n^2)$ per $z$
(or $O(n)$ with per-row envelopes), so the combined cheap floor costs the same
order as either alone.

\paragraph{Complementary looseness.}
The two cheap floors overshoot for \emph{different} reasons.  The triangular gap
is loose where the Schur off-diagonal tails $\tau_i = \sum_{j>i}
\abs{\tT_{ij}}$ are large (Lemma~\ref{lem:sandwich}); the SVD-frame bound is loose
where the singular-vector misalignment $\tW - I$ is large
(Remark~\ref{rem:svdnonnormal}).  These are distinct signatures of
non-normality, so the maximum \eqref{eq:cheapfloor} never decreases the floor and
contracts the inconclusive region.  In practice the triangular gap --- with
condition number one and the local sandwich of Lemma~\ref{lem:sandwich} ---
dominates over most of the plane, while the SVD frame contributes its sharpest
near defective or clustered spectra, where singular-vector misalignment is the
more informative measure; the relative frequency is a per-matrix empirical
question.

\paragraph{Routing consequence.}
In the adaptive sweep the cheap layer becomes ``evaluate $g_{\mathrm{tri}}$ and
$\gamma^{1/2}$, take the max''; escalation to the recentred
(Theorem~\ref{thm:resolvent-row}) or cluster (Theorem~\ref{thm:assembly})
certificates is triggered only where this combined floor remains below
threshold.  The benchmark records, per cell, which floor was binding ---
triangular, SVD-frame, or neither (escalated) --- since that frequency directly
measures whether a matrix's non-normality manifests more in Schur tails or in
singular-vector misalignment.

\end{document}
